\documentclass[12pt]{article}

\usepackage[T1]{fontenc}
\usepackage[utf8]{inputenc}
\usepackage{lmodern}
\usepackage[margin=1in]{geometry}
\usepackage{microtype}
\usepackage{amsmath,amssymb,amsthm,mathtools}
\usepackage{array,booktabs,enumitem}
\usepackage[colorlinks=true,linkcolor=blue,citecolor=blue,urlcolor=blue]{hyperref}
\usepackage[capitalize,nameinlink,noabbrev]{cleveref}

\newtheorem{theorem}{Theorem}[section]
\newtheorem{proposition}[theorem]{Proposition}
\newtheorem{lemma}[theorem]{Lemma}
\newtheorem{corollary}[theorem]{Corollary}
\newtheorem{condition}[theorem]{Condition}
\newtheorem{assumption}[theorem]{Assumption}
\theoremstyle{definition}
\newtheorem{definition}[theorem]{Definition}
\newtheorem{remark}[theorem]{Remark}
\newtheorem{example}[theorem]{Example}

\newcommand{\eps}{\varepsilon}
\newcommand{\R}{\mathbb{R}}
\newcommand{\N}{\mathbb{N}}
\newcommand{\E}{\mathbb{E}}
\newcommand{\Pb}{\mathbb{P}}
\newcommand{\1}{\mathbf{1}}
\newcommand{\cA}{\mathcal{A}}
\newcommand{\cB}{\mathcal{B}}
\newcommand{\cE}{\mathcal{E}}
\newcommand{\cF}{\mathcal{F}}
\newcommand{\cM}{\mathcal{M}}
\newcommand{\cN}{\mathcal{N}}
\newcommand{\cP}{\mathcal{P}}
\newcommand{\cQ}{\mathcal{Q}}
\newcommand{\cT}{\mathcal{T}}
\newcommand{\cS}{\mathcal{S}}
\newcommand{\cX}{\mathcal{X}}
\DeclareMathOperator{\conv}{conv}
\DeclareMathOperator{\supp}{supp}
\DeclareMathOperator{\spn}{sp}
\DeclareMathOperator{\Out}{Out}
\DeclareMathOperator{\BR}{BR}

\title{\textbf{Uniform Equilibrium from an Augmented Metastable Tree\\via a Direct Global Induced-MDP Comparison}}
\author{Pedro Aldighieri}
\date{April 2026}

\begin{document}
\maketitle





\begin{abstract}
Neyman's problem asks whether every finite stochastic game admits a uniform
$\eps$-equilibrium. The full conjecture remains open. This paper proves the
\emph{downstream} part of the compact-tree route under explicit conditional input:
we assume that a finite universal $\eta$-tree representation is available,
together with its first-exit controller realization, and we show that this input
yields a public finite-memory uniform $\eps$-equilibrium under a robust
nodewise irreducibility hypothesis.

Version~26 keeps the finite universal $\eta$-tree, the local Bellman
certificates $(g_{i,C}^*,h_{i,C}^*)$, the stock variable $Z$, the potential
$\Psi$, and the Kakutani fixed-point construction. The final regret estimate,
however, no longer compares local gains across different routed nodes. Once the
other players are frozen at $(\sigma_{\eta,\delta,M_*})_{-i}$, player $i$ faces a
finite one-player MDP on the finite public-memory state space of the realized
profile. Under condition~\textnormal{(a)}, this induced MDP is unichain.

Standard average-reward MDP theory therefore supplies a unique optimal gain
$g_i^*$ and an optimal bias $h_i^*$. The Bellman optimality inequality gives,
for every unilateral deviation $\tau_i$,
\[
\gamma_i^T\bigl(s,(\tau_i,(\sigma_{\eta,\delta,M_*})_{-i})\bigr)
\le
 g_i^*+\frac{\spn(h_i^*)}{T}.
\]
On the compliant side, the fixed-point construction yields an
$O(\eta+\delta)$-optimal stationary policy in the same induced MDP. If
$g_i^{\sigma}$ denotes its average reward and $\widehat h_i^{\sigma}$ a bias for
that policy, then
\[
g_i^{\sigma}\ge g_i^*-O(\eta+\delta),
\qquad
\gamma_i^T(s,\sigma_{\eta,\delta,M_*})
\ge
g_i^{\sigma}-\frac{\spn(\widehat h_i^{\sigma})}{T}.
\]
Subtracting the two estimates yields
\[
\gamma_i^T\bigl(s,(\tau_i,(\sigma_{\eta,\delta,M_*})_{-i})\bigr)
-
\gamma_i^T\bigl(s,\sigma_{\eta,\delta,M_*}\bigr)
\le
A_\eta\eta + C_\eta\delta + \frac{K_\eta}{T}.
\]
Thus the local tree machinery remains the construction engine, while the final
global comparison is the standard unichain average-reward MDP argument.
\end{abstract}






\tableofcontents

\section{Introduction}

\subsection{Neyman's problem}

Shapley introduced stochastic games in 1953 and proved existence of stationary
equilibria for the discounted evaluation \cite{Shapley1953}. The long-run average
evaluation is much subtler. For zero-sum stochastic games, Mertens and Neyman
proved existence of the uniform value \cite{MertensNeyman1981}. For two-player
non-zero-sum stochastic games, Vieille established the existence of uniform
$\eps$-equilibria \cite{Vieille2000I,Vieille2000II}. Beyond two players the full
problem remains open; see Neyman's survey chapter, Solan's chapter on the
multi-player problem, Vieille's survey, and the handbook treatment
\cite{Neyman2003,Solan2003,Vieille2002,MertensSorinZamir2003}. Positive results
are known in special subclasses, including three-player absorbing games
\cite{Solan1999}, correlated equilibrium constructions \cite{SolanVieille2002},
cyclic Markov equilibria phenomena \cite{FleschThuijsmanVrieze1997}, and extensions
to broader state-space settings \cite{Nowak2003}.

The central question is:
\begin{quote}
\textbf{Neyman's open problem.}
For every finite stochastic game and every $\eps>0$, does there exist a strategy
profile that is an $\eps$-equilibrium simultaneously for all sufficiently large
finite horizons?
\end{quote}

This paper does not claim to settle the conjecture in full generality. Its scope
is narrower and more precise: we start from the finite universal $\eta$-tree
representation produced by the route-D' reduction, together with its first-exit
controller realization, and show that once that input is available the remainder
of the route is correct under a robust nodewise irreducibility condition. The
point of the rewrite is not to alter the theorem's scope but to make every
definition used downstream mathematically correct and every stated proof complete.

\subsection{Main theorem}

We first record the standing scope. The compact asymptotic tree representation and
its first-exit controller realization are inputs to the present paper.

\begin{assumption}[Standing tree input]
\label{ass:tree-input}
For every $\eta\in(0,1)$, the game is endowed with a finite universal
$\eta$-tree representation as specified in
\Cref{ass:node-data,ass:controller-input}. All theorems in
Sections~\ref{sec:bookkeeping}--\ref{sec:proof-main} are proved relative to that
finite-dimensional input.
\end{assumption}

Under this standing input the paper proves the following theorem.

\begin{theorem}[Main theorem]
\label{thm:main}
Let
\[
G=(N,S,(A_i)_{i\in N},u,P)
\]
be a finite stochastic game, and assume \Cref{ass:tree-input}. Assume further
that condition~\textnormal{(a)} from \Cref{cond:a} holds for the augmented node
kernels of the universal $\eta$-trees. Then for every $\eps>0$ there exist a
public finite-memory strategy profile $\sigma^\eps$ and $T_0\in\N$ such that for
every $T\ge T_0$, every initial state $s\in S$, every player $i\in N$, and every
unilateral deviation $\tau_i$,
\[
\gamma_i^T(s,\sigma^\eps)
\ge
\gamma_i^T\bigl(s,(\tau_i,\sigma^\eps_{-i})\bigr)-\eps.
\]
In particular, every finite stochastic game equipped with a universal
$\eta$-tree representation and first-exit controller realization satisfying
condition~\textnormal{(a)} admits a public finite-memory uniform
$\eps$-equilibrium for every $\eps>0$.
\end{theorem}

\begin{remark}
The theorem is intentionally scoped to the downstream part of the compact-tree
route. The full Neyman conjecture remains open because the route still requires,
in the background, construction of the universal $\eta$-tree representation and
its realization map for an arbitrary game. The purpose of the present rewrite is
to show that the remaining argument from that representation to a uniform
equilibrium is mathematically correct once the earlier blockers and the
expectation-level bookkeeping repair are in place.
\end{remark}




\subsection{The v26 repair: a direct global MDP comparison}

Version~26 keeps the v10--v25 frozen-continuation, local-Nash,
bias-normalization, augmented-kernel, routing, terminal-node,
expectation-bookkeeping, finite-memory, one-profile Bellman, stock-variable,
and stagewise-$\Psi$ repairs. The only change is the global comparison step.
Pass~114 confirmed that the single-player MDP reframing is correct, but it also
pinpointed the remaining defect: the local gain $g_{i,C}^*$ is a restarted-node
average and need not equal the value of the global deviator MDP.

The repaired final step is therefore simpler and fully standard.
\begin{enumerate}[label=\textnormal{(\arabic*)},leftmargin=2.4em]
\item Freeze the other players at $(\sigma_{\eta,\delta,M_*})_{-i}$. Player $i$
faces a finite one-player MDP on the finite public-memory state space of the
realized profile. The local tree, Bellman, routing, and controller machinery are
still used to build that profile.

\item Under condition~\textnormal{(a)}, the induced MDP is unichain. Hence there
exists a unique optimal gain $g_i^*$ and an optimal bias $h_i^*$ satisfying the
average-reward Bellman optimality equation.

\item The global optimal bias gives the deviator upper bound
\[
\gamma_i^T\bigl(s,(\tau_i,(\sigma_{\eta,\delta,M_*})_{-i})\bigr)
\le
 g_i^*+\frac{\spn(h_i^*)}{T}.
\]
This is the standard finite-horizon optimality inequality for a unichain
average-reward MDP.

\item The fixed-point construction from
Sections~\ref{sec:global}--\ref{sec:realization} still supplies the compliant
stationary policy. Its global average reward is $O(\eta+\delta)$-optimal in the
same induced MDP. If $g_i^{\sigma}$ denotes that gain and
$\widehat h_i^{\sigma}$ is a bias for the compliant stationary policy, then
\[
\gamma_i^T\bigl(s,\sigma_{\eta,\delta,M_*}\bigr)
\ge
 g_i^{\sigma}-\frac{\spn(\widehat h_i^{\sigma})}{T}
\ge
 g_i^* - A_\eta^{\mathrm{opt}}\eta - C_\eta\delta - \frac{\spn(\widehat h_i^{\sigma})}{T}.
\]

\item Therefore
\[
\gamma_i^T\bigl(s,(\tau_i,(\sigma_{\eta,\delta,M_*})_{-i})\bigr)
-
\gamma_i^T\bigl(s,\sigma_{\eta,\delta,M_*}\bigr)
\le
A_\eta^{\mathrm{opt}}\eta + C_\eta\delta + O\!\left(\frac1T\right).
\]
Choosing $\eta$ and $\delta$ small and then $T$ large yields the uniform
$\eps$-equilibrium.
\end{enumerate}

No routed comparison of the family $(g_{i,C}^*)_{C\in\cT_\eta}$ along different
node sequences is used. The local gains remain essential in the construction,
but the final regret bound is now global and MDP-theoretic.

\subsection{Roadmap}



Section~\ref{sec:model} recalls the game model and states the finite-dimensional
node data used by the compact-tree route. Section~\ref{sec:bookkeeping}
introduces the bookkeeping kernel and proves the graph theorem for the internal
occupation package. Section~\ref{sec:local} formulates the local game on the
genuine augmented state space and establishes local best-response and local
equilibrium existence. Section~\ref{sec:global} defines the global correspondence
with frozen continuation values and proves a Kakutani fixed point theorem.
Section~\ref{sec:realization} defines actual node visits and realizes the fixed
point by public finite-memory visit controllers. Section~\ref{sec:proof-main}
proves the augmented-kernel bias bound, the deviation-value Lipschitz lemma, and
the stagewise potential drift lemma, and then replaces the failed routed
comparison by a direct comparison in the global induced MDP. The appendices
record the finite-MDP facts, including the unichain optimality proposition and
the induced-MDP comparison used in the proof of \Cref{thm:main}, together with
the stagewise summation lemma.

\section{Model and standing asymptotic input}
\label{sec:model}

\subsection{Finite stochastic games}

A finite stochastic game is a tuple
\[
G=(N,S,(A_i)_{i\in N},u,P)
\]
where $N=\{1,\dots,n\}$ is the finite set of players, $S$ is the finite state
space, $A_i(s)$ is the finite action set of player $i$ at state $s$, the stage
payoff of player $i$ is a function
\[
u_i:S\times \prod_{j\in N}A_j(s)\to [0,1],
\]
and the transition kernel is
\[
P(\cdot\mid s,a)\in \Delta(S)
\qquad (s\in S,\ a\in \prod_{j\in N}A_j(s)).
\]
The public history at stage $t$ is
\[
h_t=(s_1,a_1,s_2,a_2,\dots,s_t).
\]
A behavioral strategy of player $i$ assigns to every public history $h_t$ a
probability distribution on $A_i(s_t)$.

For a horizon $T\ge 1$, an initial state $s$, and a profile $\sigma$, the average
payoff of player $i$ is
\[
\gamma_i^T(s,\sigma)
=
\E_{s,\sigma}\Bigl[\frac1T\sum_{t=1}^T u_i(s_t,a_t)\Bigr].
\]
We write
\[
\|u\|_\infty
:=
\max_{i\in N}\max_{s\in S}\max_{a\in\prod_{j\in N}A_j(s)} |u_i(s,a)|.
\]
Under the standing normalization $u_i\in[0,1]$, one has $\|u\|_\infty\le 1$.

\begin{definition}
A profile $\sigma$ is an \emph{$\eps$-equilibrium for horizon $T$} if for every
player $i$, every initial state $s$, and every unilateral deviation $\tau_i$,
\[
\gamma_i^T(s,\sigma)
\ge
\gamma_i^T\bigl(s,(\tau_i,\sigma_{-i})\bigr)-\eps.
\]
It is a \emph{uniform $\eps$-equilibrium} if there exists $T_0$ such that the same
inequality holds for every $T\ge T_0$.
\end{definition}

\subsection{Node data from the tree representation}

The present paper starts after the route-D' reduction has produced a finite rooted
universal $\eta$-tree. We record exactly the node data used later.

\begin{assumption}[Finite node data]
\label{ass:node-data}
Fix $\eta\in(0,1)$. The universal $\eta$-tree $\cT_\eta$ is finite. We write
$C_{\mathrm{rt}}$ for its unique root. Its nodes are of two kinds: ordinary nodes
and degenerate terminal nodes. Every node $C\in \cT_\eta$ carries the following
data.
\begin{enumerate}[label=\textnormal{(D\arabic*)},leftmargin=2.8em]
\item A finite phase-state set
\[
\Xi_C=\bigsqcup_{k\in\mathbb Z/d_C\mathbb Z}\{k\}\times S_C^k
\]
and a distinguished reference state $\xi_C^\circ\in\Xi_C$.
\item For each player $i$ and phase-state $\xi\in\Xi_C$, a finite action set
$A_i(\xi)$.
\item A finite boundary set $B_C$ and a successor map $b\mapsto C[b]$ sending
boundary states to child nodes of $\cT_\eta$. Whenever the earlier route would
have attached a terminal label $b$, the enlarged tree now contains the
corresponding degenerate terminal child $T_b:=C[b]$.
\item Internal transition coefficients
\[
P_C^0(\xi'\mid \xi,a)\in[0,1],
\qquad \xi,\xi'\in \Xi_C,
\ a\in \prod_{j\in N}A_j(\xi),
\]
and exit coefficients
\[
\lambda_C(b\mid \xi,a)\in[0,1],
\qquad b\in B_C,
\]
such that
\[
\sum_{\xi'\in\Xi_C}P_C^0(\xi'\mid \xi,a)+\sum_{b\in B_C}\lambda_C(b\mid\xi,a)=1.
\]
\item A bookkeeping restart kernel fixed by the convention
\[
\rho_C(\cdot\mid b)=\delta_{\xi_C^\circ}
\qquad (b\in B_C).
\]
This kernel is used only in the internal bookkeeping chain from
Section~\ref{sec:bookkeeping}; it is not part of the executed first-exit
dynamics. In particular every bookkeeping restart returns to the reference
state $\xi_C^\circ$.
\item A universal bias box constant $B_\eta^0<\infty$ such that every local
Bellman certificate used later can be normalized inside
$[-B_\eta^0,B_\eta^0]^{\Xi_C}$.
\end{enumerate}

If $C=T$ is a degenerate terminal node, then
\[
\Xi_T=\{\xi_T^\circ\},
\qquad
B_T=\varnothing,
\qquad
A_i(\xi_T^\circ)\text{ is a singleton for every }i,
\]
the unique internal transition is absorbing,
\[
P_T^0(\xi_T^\circ\mid \xi_T^\circ,a)=1,
\]
and for the unique action profile $a^\dagger$ one has
\[
u_i(\xi_T^\circ,a^\dagger)=w_i(T)\in[0,1].
\]
Thus $T$ is a one-state absorbing leaf with trivial bookkeeping kernel and exact
local certificate
\[
g_{i,T}=w_i(T),
\qquad
h_{i,T}\equiv 0.
\]
For a boundary label $b$ leading to $T_b$, we also write $w(b):=w(T_b)$.
\end{assumption}

\begin{remark}[Tree-depth bookkeeping]
\label{rem:tree-depth}
For each degenerate terminal node $T$, set $d_\eta(T):=1$. For an ordinary node
$C\in\cT_\eta$, define recursively
\[
d_\eta(C)
:=
\begin{cases}
1, & B_C=\varnothing,\\[0.3em]
1+\max_{b\in B_C} d_\eta(C[b]), & B_C\neq \varnothing.
\end{cases}
\]
Thus $d_\eta(C)$ is the maximal number of node layers on a path beginning at $C$.
Set
\[
D_\eta:=\max_{C\in\cT_\eta} d_\eta(C)<\infty.
\]
Whenever $b\in B_C$, the remaining depth below $b$ is at most $d_\eta(C)-1$.
Accordingly, boundary continuation values will always be stored in
$[0,d_\eta(C)-1]^N$, and the corresponding local gains will lie in
$[0,d_\eta(C)]$. In particular a exit-routing path contains at most
$D_\eta-1$ strict child moves, hence certainly at most $D_\eta$ strict exit moves.
Once play enters a degenerate terminal node, no further strict exit can occur.
\end{remark}


\begin{assumption}[First-exit controller realization]
\label{ass:controller-input}
Fix $\eta\in(0,1)$ and a node $C\in\cT_\eta$. For every stationary mixed kernel
$\sigma_C\in\Sigma_C$ and every cap $L\ge 1$, the universal
$\eta$-tree representation supplies a public finite-memory first-exit
\emph{visit controller} $\mathcal R_{C,L}^{\sigma_C}$. During every live stage
with current phase-state $\xi\in\Xi_C$, the compliant action profile is sampled
from the prescribed stationary kernel $\sigma_C(\cdot\mid\xi)$; the additional
public memory is only an implementation device and does not alter the one-step
coefficients on $\Xi_C$. The controller has the following properties.
\begin{enumerate}[label=\textnormal{(R\arabic*)},leftmargin=2.8em]
\item If the controller starts from an entry phase-state in $\Xi_C$, its path
evolves on the augmented state space
\[
\Xi_C^{\mathrm{aug}}=\Xi_C\sqcup\partial B_C
\]
until the stopping time
\[
\tau_{C,L}
:=
\inf\{t\ge 1:X_t\in\partial B_C\}\wedge L.
\]
This stopping time is the length of the actual node visit. If
$\tau_{C,L}<L$ and $X_{\tau_{C,L}}=\partial_b$, then the current visit ends at
that stage and the public routing automaton launches the successor node $C[b]$
at the next stage. No payoff is attached to the abstract exit state
$\partial_b$ itself.
\item At most one boundary label is emitted during a visit. If
$\tau_{C,L}<L$ and $X_{\tau_{C,L}}=\partial_b$, then that unique label is $b$.
If $\tau_{C,L}=L$ without a boundary hit, the visit ends by cap and the next
visit starts in the same node.
\item Let $N_{C,L}(\xi,a)$ be the number of live stages
$t\in\{0,\dots,\tau_{C,L}-1\}$ at which $(X_t,A_t)=(\xi,a)$, and let
$I_{C,L}(b):=\1_{\{X_{\tau_{C,L}}=\partial_b\}}$. Define the empirical visit
rates
\[
\widehat\mu_{C,L}^{\mathrm{rate}}(\xi,a):=\frac{N_{C,L}(\xi,a)}{\tau_{C,L}},
\qquad
\widehat\beta_{C,L}(b):=\frac{I_{C,L}(b)}{\tau_{C,L}}.
\]
Then, uniformly over entry states, there exists $K_C(\sigma_C)<\infty$ such that
\begin{equation}
\label{eq:controller-input-occ}
\E\Bigl[
\bigl\|(\widehat\mu_{C,L}^{\mathrm{rate}},\widehat\beta_{C,L})
-
\bar\mu_C^\sigma
\bigr\|_1
\Bigr]
\le \frac{K_C(\sigma_C)}{\sqrt L}.
\end{equation}
\item The controller realizes the bookkeeping occupation package in expectation.
If $\bar\mu_C^\sigma=(\mu_C^\sigma,\beta_C^\sigma)$ is the occupation package
from Section~\ref{sec:bookkeeping}, then for every entry state, every cap
$L\ge 1$, every phase-state $\xi\in\Xi_C$, every admissible action profile
$a\in \prod_{j\in N}A_j(\xi)$, and every boundary label $b\in B_C$,
\begin{equation}
\label{eq:controller-input-expected-counts}
\E[N_{C,L}(\xi,a)]
=
\mu_C^\sigma(\xi,a)\,\E[\tau_{C,L}],
\qquad
\E[I_{C,L}(b)]
=
\beta_C^\sigma(b)\,\E[\tau_{C,L}].
\end{equation}
\end{enumerate}
\end{assumption}

\begin{remark}
If $C=T$ is a degenerate terminal node, then the controller
$\mathcal R_{T,L}^{\sigma_T}$ is the trivial capped visit of length $L$: it
stays at $\xi_T^\circ$ throughout, emits no boundary label, satisfies
$\tau_{T,L}=L$ almost surely, and yields the exact actual payoff
\[
U_{i,T,L}^{\mathrm{act}}=L\,w_i(T).
\]
Hence both \eqref{eq:controller-input-occ} and
\eqref{eq:controller-input-expected-counts} hold with zero error at degenerate
terminal nodes.
\end{remark}


\subsection{Pure and mixed stationary kernels inside a node}

Fix a node $C$. A \emph{pure stationary selector} of player $i$ in node $C$ is a
map
\[
f_i: \Xi_C\to \bigsqcup_{\xi\in\Xi_C}A_i(\xi)
\quad\text{with }f_i(\xi)\in A_i(\xi).
\]
The finite set of such selectors is denoted by $F_{i,C}$. A \emph{stationary
mixed kernel} of player $i$ is an element of
\[
\Sigma_{i,C}:=\prod_{\xi\in\Xi_C}\Delta(A_i(\xi)).
\]
We write $\Sigma_C:=\prod_{i\in N}\Sigma_{i,C}$.

\begin{remark}
A point of $\Sigma_{i,C}$ may be viewed either as a statewise mixed-action kernel
or, equivalently, as a product distribution over the finite selector set
$F_{i,C}$. We work with statewise mixed kernels because the deviation MDP
coefficients are affine in those kernels; this is the precise form needed later.
The local equilibrium existence result in Section~\ref{sec:local} is equivalent to
Nash's theorem for the finite selector game.
\end{remark}

\section{The bookkeeping node system and the linear exit map}
\label{sec:bookkeeping}

\subsection{Bookkeeping kernel and balance equations}

Fix a node $C\in\cT_\eta$. This section does \emph{not} describe the executed
block dynamics. It records only the internal bookkeeping chain used to define the
convex occupation package. Exit mass is instantaneously re-injected into internal
states through $\rho_C$, so the bookkeeping chain lives on $\Xi_C$ alone.

For a stationary mixed kernel $\sigma_C\in\Sigma_C$, define the induced internal
kernel
\[
Q_C^{\sigma}(\xi'\mid \xi)
:=
\sum_{a\in\prod_j A_j(\xi)} \sigma_C(a\mid \xi)
P_C^0(\xi'\mid \xi,a),
\]
and the induced exit kernel
\[
\Lambda_C^{\sigma}(b\mid\xi)
:=
\sum_{a\in\prod_j A_j(\xi)}\sigma_C(a\mid\xi)\lambda_C(b\mid\xi,a).
\]
The corresponding bookkeeping kernel on $\Xi_C$ is
\begin{equation}
\label{eq:bookkeeping-kernel}
\widetilde Q_C^{\sigma}(\xi'\mid \xi)
=
Q_C^{\sigma}(\xi'\mid \xi)
+
\sum_{b\in B_C}\Lambda_C^{\sigma}(b\mid \xi)\rho_C(\xi'\mid b).
\end{equation}
Its row sums are equal to one by construction.

If $\pi_C^{\sigma,\mathrm{book}}$ denotes an invariant law of
$\widetilde Q_C^{\sigma}$, define the internal occupation measure and exit rate by
\[
\mu_C^{\sigma}(\xi,a):=\pi_C^{\sigma,\mathrm{book}}(\xi)\sigma_C(a\mid\xi),
\qquad
\beta_C^{\sigma}(b):=\sum_{(\xi,a)\in\Omega_C}
\mu_C^{\sigma}(\xi,a)\lambda_C(b\mid\xi,a).
\]
The pair is denoted by
\[
\bar\mu_C^{\sigma}=(\mu_C^{\sigma},\beta_C^{\sigma}).
\]


\begin{definition}
\label{def:PC}
The \emph{internal feasible set} of node $C$ is the set of all
$\mu\in\R_+^{\Omega_C}$ for which there exists $\beta\in\R_+^{B_C}$ such that
\begin{align}
\sum_{(\xi,a)\in\Omega_C}\mu(\xi,a)&=1,
\label{eq:norm-PC}
\\
\pi_\mu(\xi')
&=
\sum_{(\xi,a)\in\Omega_C}\mu(\xi,a)P_C^0(\xi'\mid\xi,a)
+
\sum_{b\in B_C}\beta(b)\rho_C(\xi'\mid b)
\qquad (\xi'\in\Xi_C),
\label{eq:internal-balance}
\\
\beta(b)
&=
\sum_{(\xi,a)\in\Omega_C}\mu(\xi,a)\lambda_C(b\mid\xi,a)
\qquad (b\in B_C),
\label{eq:beta-balance}
\end{align}
where
\[
\pi_\mu(\xi):=\sum_{a\in\prod_j A_j(\xi)}\mu(\xi,a).
\]
We denote this set by $\cP_C$.
\end{definition}


\begin{remark}
The restart kernel $\rho_C$ belongs only to the bookkeeping balance equations
\eqref{eq:internal-balance}. It does not describe what the executed dynamics do
after leaving node $C$. In the realized strategy a node visit stops at the first
exit state, and the public routing automaton moves immediately to the successor
node at the next stage; the restart kernel never appears in the executed play.
\end{remark}

\subsection{The graph theorem for exit fluxes}

Define the linear exit map
\begin{equation}
\label{eq:LC-def}
L_C:\R^{\Omega_C}\to \R^{B_C},
\qquad
(L_C\mu)(b):=
\sum_{(\xi,a)\in\Omega_C}\mu(\xi,a)\lambda_C(b\mid\xi,a).
\end{equation}
The associated restart operator is
\begin{equation}
\label{eq:RC-def}
R_C:\R^{B_C}\to \R^{\Xi_C},
\qquad
(R_C\beta)(\xi'):=\sum_{b\in B_C}\beta(b)\rho_C(\xi'\mid b).
\end{equation}


\begin{proposition}
\label{prop:PC-polytope}
For every node $C$, the set $\cP_C$ is a nonempty compact convex polytope.
Moreover, $\mu\in\cP_C$ if and only if
\begin{align}
\sum_{(\xi,a)}\mu(\xi,a)&=1,
\label{eq:PC-short1}
\\
\pi_\mu &= P_C^0\mu + R_C L_C\mu,
\label{eq:PC-short2}
\end{align}
where $(P_C^0\mu)(\xi'):=\sum_{(\xi,a)}\mu(\xi,a)P_C^0(\xi'\mid\xi,a)$.
\end{proposition}

\begin{proof}
The defining relations in \Cref{def:PC} are finitely many linear equalities and
inequalities in the variables $(\mu,\beta)$. Hence the feasible set in
$\R_+^{\Omega_C}\times\R_+^{B_C}$ is a compact convex polyhedron; compactness
comes from \eqref{eq:norm-PC}. Projecting onto the $\mu$-coordinates shows that
$\cP_C$ is a compact convex polytope.

If $\mu\in\cP_C$, the auxiliary vector $\beta$ is unique by
\eqref{eq:beta-balance}, namely $\beta=L_C\mu$. Substituting this into
\eqref{eq:norm-PC} and \eqref{eq:internal-balance} gives
\eqref{eq:PC-short1}--\eqref{eq:PC-short2}. Conversely, if
\eqref{eq:PC-short1}--\eqref{eq:PC-short2} hold and one sets $\beta=L_C\mu$, then
all three relations from \Cref{def:PC} are satisfied.
\end{proof}


\begin{definition}
\label{def:MCbook}
The \emph{bookkeeping augmented feasible set} of node $C$ is
\[
\cM_C^{\mathrm{book}}
:=
\{(\mu,\beta)\in\cP_C\times\R_+^{B_C}:\beta=L_C\mu\}.
\]
The affine graph map is
\[
J_C:\cP_C\to \cM_C^{\mathrm{book}},
\qquad J_C(\mu)=(\mu,L_C\mu).
\]
\end{definition}

\begin{theorem}[Graph theorem for the bookkeeping node]
\label{thm:graph-theorem}
For every node $C$,
\[
\cM_C^{\mathrm{book}} = J_C(\cP_C).
\]
The map $J_C$ is affine and injective. In particular the exit coordinates are not
independent variables: once the internal part $\mu$ is fixed, the exit mass is
uniquely determined as $L_C\mu$.
\end{theorem}

\begin{proof}
The identity is immediate from \Cref{def:MCbook}. Affineness follows from the
linearity of $L_C$, and injectivity is obvious from the first coordinate.
\end{proof}

\begin{corollary}
\label{cor:no-vp3}
There is no separate compatibility condition between internal occupation and exit
mass. The feasible bookkeeping set is exactly the graph of the linear exit map.
\end{corollary}

\begin{proof}
This is a restatement of \Cref{thm:graph-theorem}.
\end{proof}

\subsection{Actual bookkeeping packages induced by stationary mixed kernels}

\begin{definition}
\label{def:actual-local-package}
For a stationary mixed kernel $\sigma_C\in\Sigma_C$, its \emph{actual bookkeeping
occupation} is the pair $\bar\mu_C^{\sigma}$ induced by an invariant law of the
bookkeeping kernel \eqref{eq:bookkeeping-kernel}. By construction,
$\bar\mu_C^{\sigma}\in \cM_C^{\mathrm{book}}$.
\end{definition}

\begin{proof}
Let $\pi=\pi_C^{\sigma,\mathrm{book}}$ be an invariant law of
$\widetilde Q_C^{\sigma}$. Writing the stationarity equation for
$\widetilde Q_C^{\sigma}$ and expanding \eqref{eq:bookkeeping-kernel} gives
exactly \eqref{eq:norm-PC}--\eqref{eq:beta-balance} with
$\mu=\mu_C^{\sigma}$ and $\beta=\beta_C^{\sigma}$.
\end{proof}

\section{Local average-reward games and Bellman certificates}
\label{sec:local}

\subsection{Local continuation values}

Fix $\eta\in(0,1)$ and a node $C$. Let $x$ be a package of continuation data on the
whole tree. For each boundary state $b\in B_C$ the successor $D=C[b]$ is always a
child node of the enlarged tree. The value
$v_C^x(b)\in[0,d_\eta(C)-1]^N$ is the continuation vector stored at that boundary
coordinate in $x$. At a fixed point it is tied, via
\eqref{eq:frozen-continuation}, to the child-node continuation benchmark, and in
the global proof below it is exactly the datum that transfers the Bellman
potential across a strict exit. If $D=T_b$ is a degenerate terminal child, then
$g_D^x=w(D)=w(b)$ exactly and $h_D^x=0$. Thus exits to degenerate terminal nodes
enter the local Bellman equations through the same frozen continuation
coordinates as all other exits. There is no separate terminal accounting outside
the fixed-point system. Throughout this section $x$ is considered fixed.


\subsection{The stopped kernel and the continuation Bellman operator}

Fix a node $C$ and a stationary mixed kernel $\sigma_C\in\Sigma_C$. The
\emph{augmented state space} is
\[
\Xi_C^{\mathrm{aug}}:=\Xi_C\sqcup \partial B_C,
\qquad
\partial B_C:=\{\partial_b:b\in B_C\}.
\]
The executed kernel of a capped node visit is
\begin{equation}
\label{eq:aug-kernel}
P_C^{\sigma,\mathrm{aug}}(z\mid y)
=
\begin{cases}
Q_C^{\sigma}(\xi'\mid \xi), & y=\xi\in\Xi_C,\ z=\xi'\in\Xi_C,\\[0.3em]
\Lambda_C^{\sigma}(b\mid \xi), & y=\xi\in\Xi_C,\ z=\partial_b,\\[0.3em]
1, & y=z=\partial_b,\\[0.3em]
0, & \text{otherwise.}
\end{cases}
\end{equation}
The absorbing exit states $\partial_b$ are a stopping-time device: they record
that the current node visit has ended through boundary label $b$. They are not
stages of the stochastic game and they do not carry any literal stage payoff.

For a fixed player $i$ and opponents' kernel $\sigma_{-i}$, the stage reward part,
the internal transition part, and the exit part are
\begin{align}
\label{eq:r-sigma}
r_{i,C}^{\sigma_{-i}}(\xi,a_i)
&:=
\sum_{a_{-i}}\sigma_{-i}(a_{-i}\mid\xi)
\,u_i\bigl(\xi,(a_i,a_{-i})\bigr),
\\
\label{eq:Q-sigma}
Q_{i,C}^{\sigma_{-i}}(\xi'\mid \xi,a_i)
&:=
\sum_{a_{-i}}\sigma_{-i}(a_{-i}\mid\xi)
\,P_C^0\bigl(\xi'\mid\xi,(a_i,a_{-i})\bigr),
\\
\label{eq:q-exit-sigma}
\Lambda_{i,C}^{\sigma_{-i}}(b\mid\xi,a_i)
&:=
\sum_{a_{-i}}\sigma_{-i}(a_{-i}\mid\xi)
\,\lambda_C\bigl(b\mid\xi,(a_i,a_{-i})\bigr).
\end{align}
The auxiliary continuation-corrected coefficient is
\begin{equation}
\label{eq:augmented-reward}
\widetilde r_{i,C}^{\sigma_{-i},x}(\xi,a_i)
:=
r_{i,C}^{\sigma_{-i}}(\xi,a_i)
+
\sum_{b\in B_C}\Lambda_{i,C}^{\sigma_{-i}}(b\mid\xi,a_i)v_{i,C}^x(b).
\end{equation}
This is \emph{not} an executed stage payoff. It is the coefficient appearing in
the Bellman equation for the stopped node-visit problem: one sums the actual
stage payoffs earned before the visit stops and then attaches a single
continuation certificate at the boundary label, if a boundary is hit.

The Bellman operator is therefore
\begin{equation}
\label{eq:Bellman-operator}
(T_{i,C}^{\sigma_{-i},x}h)(\xi)
:=
\max_{a_i\in A_i(\xi)}\Bigl[
 r_{i,C}^{\sigma_{-i}}(\xi,a_i)
 +\sum_{\xi'\in\Xi_C}Q_{i,C}^{\sigma_{-i}}(\xi'\mid\xi,a_i)h(\xi')
 +\sum_{b\in B_C}\Lambda_{i,C}^{\sigma_{-i}}(b\mid\xi,a_i)v_{i,C}^x(b)
 \Bigr].
\end{equation}
Equivalently, this is the Bellman operator of the stopped augmented MDP on
$\Xi_C^{\mathrm{aug}}$ with boundary biases fixed at zero.

\begin{lemma}
\label{lem:coeff-affine}
For fixed $C$, $i$, and $x$, the maps
\[
\sigma_{-i}\mapsto r_{i,C}^{\sigma_{-i}},
\qquad
\sigma_{-i}\mapsto Q_{i,C}^{\sigma_{-i}},
\qquad
\sigma_{-i}\mapsto \Lambda_{i,C}^{\sigma_{-i}},
\qquad
\sigma_{-i}\mapsto \widetilde r_{i,C}^{\sigma_{-i},x}
\]
are affine. Consequently there exists a finite constant $K_{i,C}$ such that for
all $\sigma_{-i},\tau_{-i}$,
\begin{align}
\|\widetilde r_{i,C}^{\tau_{-i},x}-\widetilde r_{i,C}^{\sigma_{-i},x}\|_\infty
&\le K_{i,C}\|\tau_{-i}-\sigma_{-i}\|_1,
\label{eq:r-lip}
\\
\sup_{\xi,a_i}\|Q_{i,C}^{\tau_{-i}}(\cdot\mid\xi,a_i)-Q_{i,C}^{\sigma_{-i}}(\cdot\mid\xi,a_i)\|_{\mathrm{TV}}
&\le K_{i,C}\|\tau_{-i}-\sigma_{-i}\|_1,
\label{eq:Q-lip}
\\
\sup_{\xi,a_i}\|\Lambda_{i,C}^{\tau_{-i}}(\cdot\mid\xi,a_i)-\Lambda_{i,C}^{\sigma_{-i}}(\cdot\mid\xi,a_i)\|_{1}
&\le K_{i,C}\|\tau_{-i}-\sigma_{-i}\|_1.
\label{eq:Lambda-lip}
\end{align}
\end{lemma}

\begin{proof}
Each coefficient is a finite linear combination of the components of
$\sigma_{-i}(\cdot\mid\xi)$, with coefficients given by the primitive reward,
transition, and continuation data. This proves affineness. The Lipschitz bounds
follow because all underlying arrays are bounded and finite.
\end{proof}

\subsection{Local best responses}

We now formulate the local best-response correspondence. Since the local state and
action sets are finite, every unilateral local problem is a finite average-reward
MDP with communicating internal part and bounded continuation rewards at exits.

\begin{definition}
\label{def:local-BR-set}
For fixed $x$, node $C$, and player $i$, let $\mathfrak B_{i,C}(x)$ be the set of
triples $(\sigma_{i,C},g_{i,C},h_{i,C})$ such that
\begin{enumerate}[label=\textnormal{(BR\arabic*)},leftmargin=2.9em]
\item $\sigma_{i,C}\in \Sigma_{i,C}$;
\item $g_{i,C}\in[0,d_\eta(C)]$ and
$h_{i,C}\in[-B_\eta^0,B_\eta^0]^{\Xi_C}$;
\item $h_{i,C}$ is normalized by
\[
h_{i,C}(\xi_C^\circ)=0;
\]
\item for every $\xi\in\Xi_C$,
\begin{equation}
\label{eq:Bellman-ineq-def}
g_{i,C}+h_{i,C}(\xi)
\ge
r_{i,C}^{\sigma_{-i,C}^x}(\xi,a_i)
+
\sum_{\xi'}Q_{i,C}^{\sigma_{-i,C}^x}(\xi'\mid\xi,a_i)h_{i,C}(\xi')
+
\sum_b \Lambda_{i,C}^{\sigma_{-i,C}^x}(b\mid\xi,a_i)v_{i,C}^x(b)
\end{equation}
for every $a_i\in A_i(\xi)$;
\item equality holds in \eqref{eq:Bellman-ineq-def} for every action $a_i$ in the
support of $\sigma_{i,C}(\cdot\mid\xi)$.
\end{enumerate}
\end{definition}

\begin{proposition}
\label{prop:local-BR}
Assume condition~\textnormal{(a)}. For every fixed $x$, node $C$, and player $i$,
the set $\mathfrak B_{i,C}(x)$ is nonempty, compact, and convex. The
correspondence $x\mapsto \mathfrak B_{i,C}(x)$ has closed graph.
\end{proposition}

\begin{proof}
Fix $x$, $C$, and $i$. The Bellman system
\eqref{eq:Bellman-ineq-def} is exactly the dual optimality system for the finite
average-reward MDP faced by player $i$ against the frozen environment encoded by
$x$; see \Cref{prop:mdp-duality}. Because the frozen continuation vector lies in
$[0,d_\eta(C)-1]^N$ and the stage payoffs lie in $[0,1]$, the augmented one-step
reward in that MDP is bounded by $d_\eta(C)$. Hence there exists an optimal
stationary mixed kernel together with a gain-bias pair satisfying
\eqref{eq:Bellman-ineq-def} and the support equalities, with
$g_{i,C}\in[0,d_\eta(C)]$. Therefore $\mathfrak B_{i,C}(x)$ is nonempty.

Compactness follows because $\Sigma_{i,C}$ is a product of simplices and the gain
and bias coordinates are confined to compact boxes. The defining equalities and
inequalities are closed.

For convexity, condition~\textnormal{(a)} implies that the internal part of the
deviation MDP is communicating on $\Xi_C$. For communicating finite MDPs, the
optimal bias is unique up to an additive constant; after the linear normalization
$h(\xi_C^\circ)=0$, the optimal bias is unique, see
\Cref{prop:bias-unique} and \cite[Chapter~8]{Puterman1994}. The same is true for
the optimal gain. Hence, once $x$ is fixed, there is a unique normalized optimal
pair $(g^*,h^*)$. A stationary mixed kernel is optimal if and only if, at every
state, it assigns mass only to actions attaining equality in the Bellman system
for $(g^*,h^*)$. The set of such kernels is a product of simplices and hence
convex. Thus $\mathfrak B_{i,C}(x)$ is convex.

For the graph property, let $x^m\to x$ and let
$(\sigma_i^m,g_i^m,h_i^m)\in \mathfrak B_{i,C}(x^m)$ converge to
$(\sigma_i,g_i,h_i)$. The coefficient arrays from \Cref{lem:coeff-affine} depend
continuously on $x$, hence passing to the limit in the Bellman equalities and
inequalities shows $(\sigma_i,g_i,h_i)\in\mathfrak B_{i,C}(x)$.
\end{proof}

\subsection{Local equilibrium existence}

\begin{proposition}[Local stationary Nash equilibrium]
\label{prop:local-equilibrium}
Assume condition~\textnormal{(a)}. Fix $x$ and a node $C$. There exists a
stationary mixed kernel $\sigma_C^*\in\Sigma_C$ such that for every player $i$,
$(\sigma_{i,C}^*,g_{i,C}^*,h_{i,C}^*)\in \mathfrak B_{i,C}(x)$ for some
$(g_{i,C}^*,h_{i,C}^*)$. Equivalently, the local average game at node $C$ with
continuation values frozen at $x$ has a stationary Nash equilibrium, and every
player admits a Bellman certificate at that equilibrium.
\end{proposition}

\begin{proof}
Fix the continuation vector $v_C^x(\cdot)$ extracted from $x$. For each player
$i$ and every opponents' kernel
$\sigma_{-i,C}\in\prod_{j\ne i}\Sigma_{j,C}$, let
\[
\BR_{i,C}^x(\sigma_{-i,C})
:=
\{\tau_{i,C}\in\Sigma_{i,C}: \exists g,h\text{ with }
(\tau_{i,C},g,h)\text{ satisfying \eqref{eq:Bellman-ineq-def} against }
\sigma_{-i,C}\text{ and }v_C^x\}.
\]
By the same MDP-duality argument as in \Cref{prop:local-BR}, each value
$\BR_{i,C}^x(\sigma_{-i,C})$ is nonempty, compact, and convex, and the graph of
the playerwise best-response correspondence is closed. Define the product
correspondence on $\Sigma_C$ by
\[
\BR_C^x(\sigma_C):=\prod_{i\in N}\BR_{i,C}^x(\sigma_{-i,C}).
\]
Its values are therefore nonempty, compact, and convex, and its graph is closed.
Since $\Sigma_C$ is a compact convex product of simplices, Kakutani's theorem
\cite{Kakutani1941} yields a fixed point $\sigma_C^*\in\BR_C^x(\sigma_C^*)$. For
each player $i$, choose a Bellman certificate $(g_{i,C}^*,h_{i,C}^*)$ associated
with the best-response relation at that fixed point. This is exactly the required
local stationary Nash package.
\end{proof}

\begin{remark}
This proposition is the correct multiplayer replacement for the v4 appendix, which
contained only one-player LP duality. One-player LP duality is still needed, but
only to identify each player's best-response set. Existence of a simultaneous
local Nash package requires the fixed-point step above.
\end{remark}

\section{The global correspondence with frozen continuation values}
\label{sec:global}

\subsection{Ambient package space}

A global package contains, at each node and for each player, a stationary mixed
kernel and a Bellman certificate, together with nodewise continuation values.
There is \emph{no} self-consistency requirement in the ambient space.

\begin{definition}
\label{def:Aeta}
For fixed $\eta\in(0,1)$, the ambient package space $\cA_\eta$ consists of all
families
\[
x=\Bigl((\sigma_{i,C},g_{i,C},h_{i,C})_{i\in N},\ (v_C(b))_{b\in B_C}\Bigr)_{C\in\cT_\eta}
\]
such that:
\begin{enumerate}[label=\textnormal{(A\arabic*)},leftmargin=2.8em]
\item $\sigma_{i,C}\in \Sigma_{i,C}$ for all $i,C$;
\item $g_{i,C}\in[0,d_\eta(C)]$ for all $i,C$;
\item $h_{i,C}\in[-B_\eta^0,B_\eta^0]^{\Xi_C}$ and
$h_{i,C}(\xi_C^\circ)=0$ for all $i,C$;
\item $v_C(b)\in[0,d_\eta(C)-1]^N$ for all $C$ and $b\in B_C$.
\end{enumerate}
\end{definition}

\begin{proposition}
\label{prop:Aeta-compact}
For every $\eta\in(0,1)$, the set $\cA_\eta$ is compact and convex.
\end{proposition}

\begin{proof}
Each $\Sigma_{i,C}$ is a product of simplices, hence compact and convex. The gain
coordinates lie in the depth-dependent intervals $[0,d_\eta(C)]$, the bias
coordinates lie in a compact box intersected with the linear normalization
hyperplane $h(\xi_C^\circ)=0$, and the continuation values lie in the
depth-dependent cubes $[0,d_\eta(C)-1]^N$. Because the tree is finite,
$\cA_\eta$ is a finite product of compact convex sets.
\end{proof}

\subsection{Global best response}

The key point is that continuation values are frozen from the reference package
$x$. The response package $y$ is allowed to differ from $x$.

\begin{definition}
\label{def:global-BR}
For $x\in\cA_\eta$, define $\BR_\eta(x)$ to be the set of all packages
$y\in\cA_\eta$ such that:
\begin{enumerate}[label=\textnormal{(G\arabic*)},leftmargin=2.8em]
\item for every node $C$ and every player $i$,
\[
(\sigma_{i,C}^y,g_{i,C}^y,h_{i,C}^y)\in \mathfrak B_{i,C}(x);
\]
\item whenever $b\in B_C$ and $D=C[b]$, one has
\begin{equation}
\label{eq:frozen-continuation}
\|v_C^y(b)-g_D^x\|_\infty\le \eta,
\end{equation}
where $g_D^x=(g_{i,D}^x)_{i\in N}$.
\end{enumerate}
\end{definition}

\begin{remark}
The response package $y$ is assembled against \emph{frozen} continuation values and
frozen opponents from $x$. There is no circularity here. One does not ask $y$ to
be self-consistent while defining $\BR_\eta(x)$; that happens only at a fixed
point. This is exactly the continuation repair already made in v5 and retained in
v6.
\end{remark}

\begin{theorem}[Global Kakutani fixed point]
\label{thm:kakutani-global}
Assume condition~\textnormal{(a)}. For every $\eta\in(0,1)$, the correspondence
$\BR_\eta:\cA_\eta\rightrightarrows\cA_\eta$ has nonempty compact convex values and
closed graph. Consequently it has a fixed point $x^*\in\cA_\eta$.
\end{theorem}

\begin{proof}
By \Cref{prop:Aeta-compact}, the domain $\cA_\eta$ is compact and convex.

Fix $x\in\cA_\eta$. For every node $C$ and player $i$, the set
$\mathfrak B_{i,C}(x)$ is nonempty, compact, and convex by
\Cref{prop:local-BR}. The continuation constraints
\eqref{eq:frozen-continuation} are closed convex constraints. Therefore
$\BR_\eta(x)$ is a finite product of nonempty compact convex sets, hence itself
nonempty, compact, and convex.

For the graph property, let $x^m\to x$ and $y^m\to y$ with
$y^m\in\BR_\eta(x^m)$ for all $m$. For every node $C$ and player $i$ we have
$(\sigma_{i,C}^{y^m},g_{i,C}^{y^m},h_{i,C}^{y^m})\in\mathfrak B_{i,C}(x^m)$. By
the closed-graph property from \Cref{prop:local-BR}, the limit belongs to
$\mathfrak B_{i,C}(x)$. Passing to the limit in the continuation inequalities gives
\eqref{eq:frozen-continuation} for $y$ against $x$. Hence $y\in\BR_\eta(x)$, so
$\BR_\eta$ has closed graph.

Closed graph plus compact values on a compact domain imply upper
hemicontinuity. Kakutani's theorem \cite{Kakutani1941} yields a fixed point.
\end{proof}

\begin{corollary}
\label{cor:local-nash-at-fixed-point}
If $x^*\in\BR_\eta(x^*)$, then for every node $C$ the stationary mixed kernel
$\sigma_C^*:=(\sigma_{i,C}^*)_{i\in N}$ is a local stationary Nash equilibrium of
the node game with continuation values $v_C^{x^*}(\cdot)$, and each player has a
Bellman certificate $(g_{i,C}^*,h_{i,C}^*)$.
\end{corollary}

\begin{proof}
At a fixed point,
$(\sigma_{i,C}^*,g_{i,C}^*,h_{i,C}^*)\in\mathfrak B_{i,C}(x^*)$ for every $i,C$.
Thus each player's stationary mixed kernel is a best response to the others inside
node $C$, against the continuation values carried by $x^*$. That is exactly the
local stationary Nash property.
\end{proof}


\section{Local controllers and global realization}
\label{sec:realization}

\subsection{The node controller, actual node visits, and auxiliary exit certificates}

Fix a fixed point $x^*\in\BR_\eta(x^*)$. For each node $C$, write
\[
\sigma_C^*:=(\sigma_{i,C}^*)_{i\in N},
\qquad
\bar\mu_C^*:=(\mu_C^{\sigma_C^*},\beta_C^{\sigma_C^*}).
\]
The local controller for node $C$ is the first-exit visit controller
$\mathcal R_{C,L}^{\sigma_C^*}$ supplied by \Cref{ass:controller-input}. A
launch with cap $L$ produces an \emph{actual node visit}: a contiguous interval
of stochastic-game stages during which the public routing automaton stays in the
node $C$, ending at the stopping time $\tau_{C,L}\le L$. If the visit exits
through boundary label $b$ before the cap, the next stage is routed immediately
to the successor node $C[b]$. There is no post-exit holding phase and no
artificial continuation payoff on the remaining stages.

For a node visit $V$ in node $C$, let $L(V)$ denote its announced cap and let
\[
\ell(V):=\tau(V)\le L(V)
\]
denote its realized length. Let $N_V(\xi,a)$ be the number of stages in $V$ at
which $(X_t,A_t)=(\xi,a)$, and let $I_V(b)$ be the indicator that $V$ exits
through label $b$. The empirical visit rates are
\[
\widehat\mu_V^{\mathrm{rate}}(\xi,a):=\frac{N_V(\xi,a)}{\ell(V)},
\qquad
\widehat\beta_V(b):=\frac{I_V(b)}{\ell(V)}.
\]
For any evaluation profile $\varsigma$ on the visit $V$, the \emph{actual visit
payoff} is
\begin{equation}
\label{eq:visit-actual-payoff}
U_i^{\mathrm{act}}(V;\varsigma)
:=
\sum_{t\in V} u_i^t(\varsigma).
\end{equation}
We also introduce the auxiliary \emph{exit-corrected visit functional}
\begin{equation}
\label{eq:visit-aux-payoff}
\widetilde U_i(V;\varsigma)
:=
U_i^{\mathrm{act}}(V;\varsigma)
+
\sum_{b\in B_C}I_V(b)\,v_{i,C}^{x^*}(b).
\end{equation}
This auxiliary quantity is retained only as a local stopped-Bellman diagnostic.
Version~26 keeps it local, and it also retains the stock-augmented stagewise
potential $\Psi_{i,t}=h_{i,C_t}^*(X_t)+Z_{i,t}$ defined later in
\Cref{sec:proof-main} as the correct transport of local Bellman certificates
across node exits. The final regret bound, however, is not obtained by
telescoping local gains across routed nodes. It is deduced directly from the
global gain-bias theory of the induced MDP.

Define the actual nodewise average payoff rate
\begin{equation}
\label{eq:alpha-node}
\alpha_{i,C}^*
:=
\sum_{(\xi,a)\in\Omega_C}\mu_C^*(\xi,a)u_i(\xi,a)\in[0,1].
\end{equation}

\begin{theorem}[Local controller lemma]
\label{thm:local-controller}
Assume \Cref{ass:controller-input}. Fix a node $C$ and the fixed-point kernel
$\sigma_C^*$. Then there exists a constant $K_C<\infty$ such that for every cap
$L\ge 1$, every entry phase-state, and every player $i$,
\begin{equation}
\label{eq:local-controller-occ}
\E\bigl[\|(\widehat\mu_{C,L}^{\mathrm{rate}},\widehat\beta_{C,L})-\bar\mu_C^*\|_1\bigr]
\le \frac{K_C}{\sqrt L},
\end{equation}
\begin{equation}
\label{eq:local-controller-actual-payoff}
\E[U_{i,C,L}^{\mathrm{act}}]
=
\alpha_{i,C}^*\,\E[\tau_{C,L}],
\end{equation}
\begin{equation}
\label{eq:local-controller-exit-count}
\E[I_{C,L}(b)]
=
\beta_C^*(b)\,\E[\tau_{C,L}]
\qquad (b\in B_C),
\end{equation}
and
\begin{equation}
\label{eq:local-controller-aux-payoff}
\E[\widetilde U_{i,C,L}]
=
g_{i,C}^*\,\E[\tau_{C,L}].
\end{equation}
Consequently, for every $\delta>0$ there exists $L_C(\delta)$ such that whenever
$L\ge L_C(\delta)$, the occupation estimate \eqref{eq:local-controller-occ} is at
most $\delta$; the expectation identities
\eqref{eq:local-controller-actual-payoff}--\eqref{eq:local-controller-aux-payoff}
hold for every cap $L$.
\end{theorem}

\begin{proof}
The occupation estimate \eqref{eq:local-controller-occ} is exactly the controller
estimate from \Cref{ass:controller-input}, with $K_C:=K_C(\sigma_C^*)$.

For the actual payoff identity, write
\[
U_{i,C,L}^{\mathrm{act}}
=
\sum_{\xi\in\Xi_C}
\sum_{a\in\prod_j A_j(\xi)}
N_{C,L}(\xi,a)\,u_i(\xi,a).
\]
Taking expectations and using \eqref{eq:controller-input-expected-counts} for the
kernel $\sigma_C^*$ gives
\[
\E[U_{i,C,L}^{\mathrm{act}}]
=
\sum_{\xi,a}
\mu_C^*(\xi,a)\,u_i(\xi,a)\,\E[\tau_{C,L}]
=
\alpha_{i,C}^*\,\E[\tau_{C,L}],
\]
which is \eqref{eq:local-controller-actual-payoff}. The exit-count identity
\eqref{eq:local-controller-exit-count} is the second part of
\eqref{eq:controller-input-expected-counts} for the same kernel.

For the auxiliary functional,
\[
\E[\widetilde U_{i,C,L}]
=
\E[U_{i,C,L}^{\mathrm{act}}]
+
\sum_{b\in B_C} v_{i,C}^{x^*}(b)\,\E[I_{C,L}(b)].
\]
Using \eqref{eq:local-controller-actual-payoff} and
\eqref{eq:local-controller-exit-count}, we obtain
\[
\E[\widetilde U_{i,C,L}]
=
\left(
\alpha_{i,C}^*
+
\sum_{b\in B_C}\beta_C^*(b)\,v_{i,C}^{x^*}(b)
\right)\E[\tau_{C,L}].
\]
Integrating the Bellman equalities for actions in the support of
$\sigma_{i,C}^*(\cdot\mid\xi)$ against the bookkeeping package
$\bar\mu_C^*=(\mu_C^*,\beta_C^*)$ and using
\eqref{eq:internal-balance}--\eqref{eq:beta-balance} yields
\begin{equation}
\label{eq:node-gain-identity}
g_{i,C}^*
=
\alpha_{i,C}^*
+
\sum_{b\in B_C}\beta_C^*(b)\,v_{i,C}^{x^*}(b)
-
\sum_{\xi'\in\Xi_C}(R_C\beta_C^*)(\xi')h_{i,C}^*(\xi').
\end{equation}
Because the bookkeeping restart is fixed at the reference state,
$\rho_C(\cdot\mid b)=\delta_{\xi_C^\circ}$, the last term is
\[
\Bigl(\sum_{b\in B_C}\beta_C^*(b)\Bigr)h_{i,C}^*(\xi_C^\circ)=0
\]
by the normalization $h_{i,C}^*(\xi_C^\circ)=0$. Therefore
\[
g_{i,C}^*
=
\alpha_{i,C}^*
+
\sum_{b\in B_C}\beta_C^*(b)\,v_{i,C}^{x^*}(b),
\]
and substituting this identity into the previous display gives
\eqref{eq:local-controller-aux-payoff}.
\end{proof}


\subsection{Frozen superexponential schedule}

\begin{definition}
A sequence $(M_r)_{r\ge 1}$ of positive integers is a
\emph{superexponential schedule} if $M_{r+1}\ge M_r^2$ for all $r$. Its
\emph{freezing at level $R$} is the sequence $(M_r^{(R)})_{r\ge 1}$ defined by
\[
M_r^{(R)}=
\begin{cases}
M_r, & r<R,\\
M_R, & r\ge R.
\end{cases}
\]
We write $M_*:=M_R$ for the tail visit cap.
\end{definition}

\begin{lemma}
\label{lem:frozen-memory}
A frozen superexponential visit schedule uses only finite public memory.
\end{lemma}

\begin{proof}
Only the finitely many caps $M_1,\dots,M_R$ appear. To implement the schedule one
stores the current node, the current position inside the active visit, the
current local memory state of the node controller, and the visit index truncated
at $R$. All of these coordinates range over finite sets.
\end{proof}

\begin{definition}[Actual node-visit schedule]
\label{def:visit-schedule}
Fix a frozen superexponential schedule and a public routing automaton. Let
$t_1:=1$ and $C_1:=C_{\mathrm{rt}}$. At each launch time $t_k$, the public memory
determines the current node $C_k$ and an announced cap
\[
L_k:=M_k^{(R)}.
\]
The automaton launches the controller
$\mathcal R_{C_k,L_k}^{\sigma_{C_k}^*}$ and defines the $k$th \emph{node visit}
by
\[
V_k:=[t_k,t_k+\ell_k),
\qquad
\ell_k:=\tau(V_k)\le L_k.
\]
The next launch time is
\[
t_{k+1}:=t_k+\ell_k.
\]
If $V_k$ exits through label $b$, then $C_{k+1}=C_k[b]$; otherwise
$C_{k+1}=C_k$.
\end{definition}

\subsection{Public routing and realization}

Recall from \Cref{rem:tree-depth} that $D_\eta<\infty$ is the maximal number of
node layers on a root-to-leaf path of the finite universal tree $\cT_\eta$.

\begin{proposition}[Public routing]
\label{prop:public-routing}
Fix a fixed point $x^*$ and a frozen superexponential schedule. There exists a
public finite-memory profile $\sigma_{\eta,\delta,M_*}$ with the following
properties.
\begin{enumerate}[label=\textnormal{(R\arabic*)},leftmargin=2.8em]
\item At each launch time $t_k$ the public routing automaton reads the current
public memory, chooses a node $C_k$ and the announced cap $L_k$ from the frozen
schedule, and launches the controller $\mathcal R_{C_k,L_k}^{\sigma_{C_k}^*}$.
The resulting actual segment is the node visit $V_k$ from
\Cref{def:visit-schedule}.
\item If $V_k$ exits through boundary label $b$ before its cap, then the next
stage is routed immediately to the successor node $D=C_k[b]$, and the next visit
starts there. If $V_k$ reaches its cap without exit, the next visit starts in the
same node.
\item Because every exit-routing move goes to a strict child and degenerate
terminal nodes have no exits, every play path contains at most $D_\eta$ strict
node changes.
\item The public memory stores the current node, the current cap index, the
current position inside the current active visit, and the local controller state.
Hence the profile uses finite public memory.
\end{enumerate}
\end{proposition}

\begin{proof}
Item~(R1) is the frozen schedule together with the controller supplied by
\Cref{ass:controller-input}. Item~(R2) is immediate from the first-exit routing
rule in that assumption. Item~(R3) follows from the rooted-tree structure: every
strict child move lowers the remaining depth and therefore can occur at most
$D_\eta$ times along a play path. Finite memory is \Cref{lem:frozen-memory}.
\end{proof}

\begin{theorem}[Realization theorem]
\label{thm:realization}
Assume condition~\textnormal{(a)} and \Cref{ass:controller-input}. Fix
$\eta\in(0,1)$ and a fixed point $x^*\in\BR_\eta(x^*)$. Let
$\sigma_{\eta,\delta,M_*}$ be the public profile obtained from
\Cref{prop:public-routing} by freezing a superexponential schedule at level $R$.
Then there exist constants $C_\eta,K_\eta<\infty$ such that:
\begin{enumerate}[label=\textnormal{(Q\arabic*)},leftmargin=2.8em]
\item if a complete tail visit $V$ is played in node $C$ and its announced cap
satisfies $L(V)\ge M_*$, then
\begin{equation}
\label{eq:visit-occupation-realization}
\E\bigl[\|(\widehat\mu_V^{\mathrm{rate}},\widehat\beta_V)-\bar\mu_C^*\|_1
\mid \cF_{\mathrm{in}(V)}\bigr]
\le \delta;
\end{equation}
\item for every player $i$ and every complete tail visit $V$ played in node $C$
with announced cap $L(V)\ge M_*$,
\begin{align}
\label{eq:visit-actual-realization-exact}
\E\bigl[U_i^{\mathrm{act}}(V;\sigma_{\eta,\delta,M_*})\mid \cF_{\mathrm{in}(V)}\bigr]
&=
\alpha_{i,C}^*\,
\E\bigl[\ell(V)\mid \cF_{\mathrm{in}(V)}\bigr],
\\
\label{eq:visit-exit-realization}
\E\bigl[I_V(b)\mid \cF_{\mathrm{in}(V)}\bigr]
&=
\beta_C^*(b)\,
\E\bigl[\ell(V)\mid \cF_{\mathrm{in}(V)}\bigr]
\qquad (b\in B_C),
\\
\label{eq:visit-aux-realization-exact}
\E\bigl[\widetilde U_i(V;\sigma_{\eta,\delta,M_*})\mid \cF_{\mathrm{in}(V)}\bigr]
&=
g_{i,C}^*\,
\E\bigl[\ell(V)\mid \cF_{\mathrm{in}(V)}\bigr],
\end{align}
and in particular
\begin{align}
\label{eq:visit-actual-realization}
\E\bigl[U_i^{\mathrm{act}}(V;\sigma_{\eta,\delta,M_*})\mid \cF_{\mathrm{in}(V)}\bigr]
&\ge
\bigl(\alpha_{i,C}^*-C_\eta\delta\bigr)\,
\E\bigl[\ell(V)\mid \cF_{\mathrm{in}(V)}\bigr],
\\
\label{eq:visit-aux-realization}
\E\bigl[\widetilde U_i(V;\sigma_{\eta,\delta,M_*})\mid \cF_{\mathrm{in}(V)}\bigr]
&\ge
\bigl(g_{i,C}^*-C_\eta\delta\bigr)\,
\E\bigl[\ell(V)\mid \cF_{\mathrm{in}(V)}\bigr];
\end{align}
\item for every profile $\varsigma$ following the same public visit-routing
convention and every horizon $T$,
\begin{equation}
\label{eq:realization-actual-tail-bound}
\left|
\gamma_i^T(s,\varsigma)
-
\frac1T\E\left[
\sum_{V\in\mathcal V_T(\varsigma)}
U_i^{\mathrm{act}}(V;\varsigma)
\right]
\right|
\le \frac{K_\eta}{T},
\end{equation}
where $\mathcal V_T(\varsigma)$ denotes the complete tail visits contained in
$[1,T]$. The constant $K_\eta$ absorbs the finite prelude visits and the final
incomplete fragment cut by the horizon;
\item the total auxiliary boundary correction is uniformly bounded:
for every profile $\varsigma$ following the public visit-routing convention and
every horizon $T$,
\begin{equation}
\label{eq:aux-actual-discrepancy}
0
\le
\sum_{V\in\mathcal V_T(\varsigma)}
\bigl(
\widetilde U_i(V;\varsigma)-U_i^{\mathrm{act}}(V;\varsigma)
\bigr)
\le
V_\eta,
\end{equation}
where one may take
\[
V_\eta:=D_\eta(D_\eta-1).
\]
\end{enumerate}
\end{theorem}

\begin{proof}
Item~(Q1) is \Cref{thm:local-controller} at the chosen tail scale, conditioned on
the entry state of the visit.

For (Q2), condition on $\cF_{\mathrm{in}(V)}$ and apply
\Cref{thm:local-controller} inside the announced node $C$ with kernel
$\sigma_C^*$. This yields the exact identities
\eqref{eq:visit-actual-realization-exact}--\eqref{eq:visit-aux-realization-exact}.
The lower bounds \eqref{eq:visit-actual-realization} and
\eqref{eq:visit-aux-realization} are immediate after enlarging $C_\eta$
uniformly over the finite tree.

For (Q3), the horizon $[1,T]$ consists of the finitely many prelude visits before
the frozen schedule reaches the tail cap $M_*$, the complete tail visits
contained in $[1,T]$, and at most one final incomplete fragment. Every prelude
visit has deterministic cap bounded by the finite prefix
$M_1,\dots,M_{R-1}$, and the final incomplete fragment has length at most $M_*$.
Therefore the total expected contribution of the prelude visits and final
fragment is bounded by a finite constant depending only on the frozen schedule
and the finite tree object. Absorb it into $K_\eta$ to obtain
\eqref{eq:realization-actual-tail-bound}.

For (Q4), each complete tail visit contributes exactly one auxiliary boundary
certificate if and only if it exits through a strict child label. Such a
certificate has size at most
\[
v_{i,C}^{x^*}(b)\le d_\eta(C)-1\le D_\eta-1.
\]
By \Cref{prop:public-routing}, there are at most $D_\eta$ strict child moves
along any play path. Hence the total auxiliary correction over the entire horizon
is bounded by $D_\eta(D_\eta-1)$.
\end{proof}


\section{Condition \texorpdfstring{$(a)$}{(a)}, bias bounds, and the proof of the main theorem}
\label{sec:proof-main}

\subsection{Condition \texorpdfstring{$(a)$}{(a)} in primitive data}

We now formulate the nodewise irreducibility hypothesis directly from the genuine
augmented transition data.

\begin{condition}[Robust nodewise irreducibility]
\label{cond:a}

For every $\eta\in(0,1)$, every node $C\in\cT_\eta$, every player $i\in N$, every
statewise mixed kernel of the other players
$\sigma_{-i}\in\prod_{j\ne i}\Sigma_{j,C}$, and every pure action selection
$a_i(\cdot)$ of player $i$, the augmented kernel
\[
P_{i,C}^{\sigma_{-i},a_i,\mathrm{aug}}(z\mid y)
\]
on $\Xi_C^{\mathrm{aug}}=\Xi_C\sqcup\partial B_C$, defined by
\[
P_{i,C}^{\sigma_{-i},a_i,\mathrm{aug}}(z\mid y)
=
\begin{cases}
\displaystyle
\sum_{a_{-i}}\sigma_{-i}(a_{-i}\mid\xi)
P_C^0\bigl(\xi'\mid\xi,(a_i(\xi),a_{-i})\bigr),
& y=\xi\in\Xi_C,\ z=\xi'\in\Xi_C,\\[0.7em]
\displaystyle
\sum_{a_{-i}}\sigma_{-i}(a_{-i}\mid\xi)
\lambda_C\bigl(b\mid\xi,(a_i(\xi),a_{-i})\bigr),
& y=\xi\in\Xi_C,\ z=\partial_b,\\[0.7em]
1, & y=z=\partial_b,\\
0, & \text{otherwise,}
\end{cases}
\]
has the property that its restriction to the internal states remains irreducible
before absorption: for every $\xi,\zeta\in\Xi_C$ there exists
$0\le n\le |\Xi_C|-1$ with
\[
\bigl(P_{i,C}^{\sigma_{-i},a_i,\mathrm{aug}}\bigr)^n(\zeta\mid\xi)>0.
\]
\end{condition}

\begin{remark}
Condition~\textnormal{(a)} has a direct operational meaning: no single player can,
by choosing actions state by state while the others keep their local kernels,
create a noncommunicating internal class inside a metastable node before the block
is absorbed at a exit state.
\end{remark}

\begin{example}
\label{ex:cond-a-trivial}
Suppose that for every node $C$, every phase-state $\xi\in\Xi_C$, and every full
action profile $a\in\prod_j A_j(\xi)$, the internal transition kernel
$P_C^0(\cdot\mid\xi,a)$ belongs to an irreducible family whose directed support
graph is the same strongly connected graph on $\Xi_C$. Then every convex mixture
of those kernels still gives positive probability to each support edge, and hence
every augmented kernel from \Cref{cond:a} can reach any internal state from any
other before absorption. Therefore condition~\textnormal{(a)} holds.
\end{example}

\subsection{A uniform irreducibility modulus}

For fixed $\eta$, define
\begin{equation}
\label{eq:kappa-eta}
\kappa_\eta
:=
\min_{C\in\cT_\eta}
\min_{i\in N}
\inf_{\sigma_{-i}\in\prod_{j\ne i}\Sigma_{j,C}}
\min_{a_i(\cdot)}
\min_{\xi,\zeta\in\Xi_C}
\max_{0\le n\le |\Xi_C|-1}
\Bigl(P_{i,C}^{\sigma_{-i},a_i,\mathrm{aug}}\Bigr)^n(\zeta\mid\xi).
\end{equation}

\begin{lemma}
\label{lem:kappa-positive}
Assume condition~\textnormal{(a)}. Then $\kappa_\eta>0$ for every $\eta\in(0,1)$.
\end{lemma}

\begin{proof}
Fix $C$, $i$, and a pure action selection $a_i(\cdot)$. By
condition~\textnormal{(a)}, every kernel
$P_{i,C}^{\sigma_{-i},a_i,\mathrm{aug}}$ has strictly positive internal
reachability between any two internal states in at most $|\Xi_C|-1$ steps. Hence
the displayed quantity in \eqref{eq:kappa-eta} is strictly positive for each fixed
$\sigma_{-i}$.

The map
$\sigma_{-i}\mapsto P_{i,C}^{\sigma_{-i},a_i,\mathrm{aug}}$ is continuous by
\Cref{lem:coeff-affine}, and the compact domain
$\prod_{j\ne i}\Sigma_{j,C}$ is a finite product of simplices. Therefore the
minimum over that domain is attained and remains strictly positive. Finally there
are only finitely many nodes, players, and pure action selections, so the overall
minimum is strictly positive.
\end{proof}

\subsection{Bias bound}


\begin{lemma}[Hitting the reference state or absorption]
\label{lem:hitting-time}
Let $X$ be a finite set, let $\partial$ be a finite absorbing set, and let $P$ be a
Markov kernel on $X\sqcup\partial$ such that for every $x,y\in X$ there exists
$0\le n\le |X|-1$ with $P^n(y\mid x)>0$. Define
\[
\kappa(P):=\min_{x,y\in X}\max_{0\le n\le |X|-1}P^n(y\mid x).
\]
Then for every $x,y\in X$,
\[
\E_x[\tau_{A_y}]\le \frac{|X|-1}{\kappa(P)},
\qquad
A_y:=\{y\}\cup \partial,
\qquad
\tau_{A_y}:=\inf\{t\ge 0:X_t\in A_y\}.
\]
\end{lemma}

\begin{proof}
If $|X|=1$, then $A_y=X\sqcup\partial$ and $\tau_{A_y}=0$, so the claim is
immediate. Assume henceforth that $|X|\ge 2$ and set $m:=|X|-1$. For each pair
$(x,y)$ choose
$n(x,y)\in\{0,\dots,m\}$ such that
\[
P^{n(x,y)}(y\mid x)\ge \kappa(P).
\]
This is possible by the definition of $\kappa(P)$. Since $n(x,y)\le m$,
\[
\Pb_x(\tau_{A_y}\le m)\ge \Pb_x(X_{n(x,y)}=y)\ge \kappa(P).
\]
Indeed the event $\{X_{n(x,y)}=y\}$ already excludes prior absorption when
$n(x,y)\ge 1$ because all states in $\partial$ are absorbing, and when
$n(x,y)=0$ it simply means $x=y$, in which case $\tau_{A_y}=0$. Therefore
\[
\Pb_x(\tau_{A_y}>m)\le 1-\kappa(P).
\]
Applying the Markov property at times $m,2m,3m,\dots$ on the event
$\{\tau_{A_y}>km\}$, where the chain is still in $X$, gives
\[
\Pb_x(\tau_{A_y}>(k+1)m)\le (1-\kappa(P))\Pb_x(\tau_{A_y}>km)
\qquad (k\ge 0),
\]
hence
\[
\Pb_x(\tau_{A_y}>km)\le (1-\kappa(P))^k.
\]
Therefore
\[
\E_x[\tau_{A_y}]
=
\sum_{t\ge 0}\Pb_x(\tau_{A_y}>t)
\le
m\sum_{k\ge 0}\Pb_x(\tau_{A_y}>km)
\le
m\sum_{k\ge 0}(1-\kappa(P))^k
=
\frac{m}{\kappa(P)}.
\]
\end{proof}


\begin{proposition}[Uniform bias bound]
\label{prop:bias-span}
Assume condition~\textnormal{(a)}. Let $x^*\in\BR_\eta(x^*)$ be a fixed point and
let $(g_{i,C}^*,h_{i,C}^*)$ be the Bellman certificate of player $i$ in node $C$,
normalized by $h_{i,C}^*(\xi_C^\circ)=0$. Then
\[
\|h_{i,C}^*\|_\infty\le
H_\eta:=D_\eta\left(1+\frac{\max_D|\Xi_D|-1}{\kappa_\eta}\right).
\]
In particular,
\[
\spn(h_{i,C}^*)\le 2H_\eta.
\]
\end{proposition}

\begin{proof}
Fix $C$ and $i$. If $C$ is degenerate terminal, then
$g_{i,C}^*=w_i(C)$ and $h_{i,C}^*\equiv 0$ by construction, so the claim is
immediate. Assume henceforth that $C$ is ordinary. Let $a_i^*(\cdot)$ be a
measurable selector choosing, in each state, an action where equality holds in
the Bellman system. Consider the genuine augmented kernel
$P_{i,C}^{\sigma_{-i,C}^*,a_i^*,\mathrm{aug}}$ on
$\Xi_C^{\mathrm{aug}}$. Let
\[
A_C^\circ:=\{\xi_C^\circ\}\cup\partial B_C,
\qquad
\tau^\circ:=\inf\{t\ge 0:X_t\in A_C^\circ\}.
\]
By \Cref{lem:kappa-positive,lem:hitting-time},
\[
\E_\xi[\tau^\circ]\le \frac{|\Xi_C|-1}{\kappa_\eta}
\qquad (\xi\in\Xi_C).
\]

Along the selected policy $a_i^*(\cdot)$, the Bellman equalities read
\[
g_{i,C}^*+h_{i,C}^*(\xi)
=
r_{i,C}^{\sigma_{-i,C}^*}(\xi,a_i^*(\xi))
+
\sum_{\xi'}Q_{i,C}^{\sigma_{-i,C}^*}(\xi'\mid\xi,a_i^*(\xi))h_{i,C}^*(\xi')
+
\sum_b\Lambda_{i,C}^{\sigma_{-i,C}^*}(b\mid\xi,a_i^*(\xi))v_{i,C}^{x^*}(b).
\]
Stopping this identity at $\tau^\circ$ and using $h_{i,C}^*(\xi_C^\circ)=0$ and
the boundary convention $h_{i,C}^*(\partial_b)=0$ gives the standard relative-value
representation
\[
h_{i,C}^*(\xi)
=
\E_\xi\!\left[
\sum_{t=0}^{\tau^\circ-1}\bigl(u_i^t-g_{i,C}^*\bigr)
+
\sum_{b\in B_C}\1_{\{X_{\tau^\circ}=\partial_b\}}\,v_{i,C}^{x^*}(b)
\right].
\]
Since $u_i^t\in[0,1]$, $g_{i,C}^*\in[0,d_\eta(C)]$, and
$v_{i,C}^{x^*}(b)\in[0,d_\eta(C)-1]$, the absolute value of the last display is
bounded by
\[
\begin{aligned}
d_\eta(C)\,\E_\xi[\tau^\circ] + d_\eta(C)-1
&\le
d_\eta(C)\bigl(\E_\xi[\tau^\circ]+1\bigr)
\\
&\le
D_\eta\left(1+\frac{|\Xi_C|-1}{\kappa_\eta}\right)
\\
&\le
D_\eta\left(1+\frac{\max_D|\Xi_D|-1}{\kappa_\eta}\right)
=
H_\eta.
\end{aligned}
\]
This proves the sup-norm bound. The span bound is immediate.
\end{proof}

\subsection{Deviation-value Lipschitz bound}

\begin{lemma}[Coefficient-level Bellman transfer]
\label{lem:coefficient-transfer}
Let $r,Q$ and $r',Q'$ be two finite average-reward MDP coefficient systems on the
same state-action graph, and let $(g,h)$ satisfy
\[
g+h(x)\ge r(x,a)+\sum_y Q(y\mid x,a)h(y)
\qquad\text{for all }(x,a).
\]
Then
\[
g+h(x)
\ge
r'(x,a)+\sum_y Q'(y\mid x,a)h(y)
-\Bigl(\|r'-r\|_\infty + \spn(h)\sup_{x,a}\|Q'(\cdot\mid x,a)-Q(\cdot\mid x,a)\|_{\mathrm{TV}}\Bigr)
\]
for all $(x,a)$. In particular the perturbation constant is governed by
$1+\spn(h)$, not merely by $\spn(h)$.
\end{lemma}

\begin{proof}
Fix $(x,a)$. Write
\begin{align*}
&r'(x,a)+\sum_y Q'(y\mid x,a)h(y)
-
\Bigl(r(x,a)+\sum_y Q(y\mid x,a)h(y)\Bigr)
\\
&= \bigl(r'(x,a)-r(x,a)\bigr)
+
\sum_y\bigl(Q'(y\mid x,a)-Q(y\mid x,a)\bigr)h(y).
\end{align*}
The first term is bounded by $\|r'-r\|_\infty$. For the second term, subtract the
minimum of $h$ and use that signed measures of total mass $0$ have dual norm equal
to total variation against functions of oscillation $\spn(h)$:
\[
\left|\sum_y\bigl(Q'-Q\bigr)(y\mid x,a)h(y)\right|
\le
\spn(h)\,\|Q'(\cdot\mid x,a)-Q(\cdot\mid x,a)\|_{\mathrm{TV}}.
\]
Combining the two bounds gives the claim.
\end{proof}

\begin{lemma}[Deviation-value Lipschitz bound]
\label{lem:deviation-lipschitz}
Fix $\eta\in(0,1)$ and a node $C$. There exists a finite constant $L_{\eta,C}$
such that for every player $i$, every two opponents' kernels
$\sigma_{-i},\tau_{-i}\in\prod_{j\ne i}\Sigma_{j,C}$, and every continuation vector
$v\in[0,d_\eta(C)-1]^{B_C\times N}$,
\[
\bigl|\operatorname{Val}_{i,C}(\tau_{-i},v)-\operatorname{Val}_{i,C}(\sigma_{-i},v)\bigr|
\le L_{\eta,C}\,\|\tau_{-i}-\sigma_{-i}\|_1,
\]
where $\operatorname{Val}_{i,C}(\cdot,v)$ denotes the optimal average reward of the
deviation MDP. One may take
\[
L_{\eta,C}=K_{i,C}(1+2H_\eta),
\]
with $K_{i,C}$ as in \Cref{lem:coeff-affine}.
\end{lemma}

\begin{proof}
Fix $\sigma_{-i}$ and let $(g,h)$ be an optimal Bellman certificate for the MDP with
coefficients induced by $(\sigma_{-i},v)$. By \Cref{prop:bias-span},
$\spn(h)\le 2H_\eta$. Apply \Cref{lem:coefficient-transfer} with the primed
coefficients induced by $(\tau_{-i},v)$. Using the Lipschitz estimates from
\Cref{lem:coeff-affine} yields
\[
\operatorname{Val}_{i,C}(\tau_{-i},v)
\le
\operatorname{Val}_{i,C}(\sigma_{-i},v)
+
K_{i,C}(1+2H_\eta)\,\|\tau_{-i}-\sigma_{-i}\|_1.
\]
Interchanging the roles of $\sigma_{-i}$ and $\tau_{-i}$ gives the reverse
inequality.
\end{proof}




\subsection{The induced one-player MDP and the stage-by-stage potential diagnostic}

The fatal v19 defect is structural. The auxiliary visit functional
\[
\widetilde U_i(V)
=
U_i^{\mathrm{act}}(V)
+
\sum_{b\in B_C}I_V(b)\,v_{i,C}^{x^*}(b)
\]
credits a strict exit, but it does not subtract the successor-node entry bias. If
a visit in node $C$ exits through $b$ and the next visit starts in $D=C[b]$, the
next visit begins with the child bias $h_{i,D}^*(X_{\mathrm{in}})$, not with
$-v_{i,C}^{x^*}(b)$. Hence no legitimate proof may telescope
$\widetilde U_i$ across visits, and the raw nodewise bias
$h_{i,C_t}^*(X_t)$ alone is not a valid global potential either.

Fix a fixed point $x^*\in\BR_\eta(x^*)$, a player $i$, and freeze the
nondeviators at the realized finite-memory profile
$(\sigma_{\eta,\delta,M_*})_{-i}$. By \Cref{prop:public-routing}, the resulting
unilateral control problem is a finite routed MDP: a state records the current
node, the current internal/controller state, and the finite public memory of the
frozen cap schedule. A unilateral deviation is exactly a policy in this MDP, and
the compliant behavior of player $i$ is the stationary policy induced by the
local kernels $\sigma_{i,C}^*$.

The Bellman packages $(g_{i,C}^*,h_{i,C}^*)$ were built precisely against this
frozen environment. The only obstruction to a global telescope is that a strict
exit from $C$ produces a Bellman continuation credit $v_{i,C}^{x^*}(b)$,
whereas the next stage starts with the child-node bias. The stock variable $Z$
is introduced exactly to transport the local Bellman certificates across those
node changes.

Set $Z_{i,1}(\varsigma):=0$; equivalently, one may view $Z_{i,0}=0$ before the
first live stage. For $t\ge 1$, define the transition stock recursively by
\begin{equation}
\label{eq:psi-stock-update}
Z_{i,t+1}(\varsigma)
=
Z_{i,t}(\varsigma)
+
\sum_{b\in B_{C_t}}
\1_{\{\textnormal{stage }t\textnormal{ exits }C_t\textnormal{ through }b\}}
\Bigl(
v_{i,C_t}^{x^*}(b)-h_{i,C_t[b]}^*(X_{t+1})
\Bigr).
\end{equation}
Thus a strict exit from $C$ to $D=C[b]$ adds exactly the Bellman exit credit
minus the successor entry bias. Define the augmented potential
\begin{equation}
\label{eq:psi-def}
\Psi_{i,t}(\varsigma)
:=
\Phi_{i,t}(\varsigma)+Z_{i,t}(\varsigma)
=
h_{i,C_t}^*(X_t)+Z_{i,t}(\varsigma).
\end{equation}
At a nonexit stage one has $Z_{i,t+1}=Z_{i,t}$, while at a strict exit through
$b$ one has
\[
\Psi_{i,t+1}(\varsigma)=Z_{i,t}(\varsigma)+v_{i,C_t}^{x^*}(b),
\]
so the jump of the raw bias is exactly compensated by the increment of $Z$.

Every stock jump is uniformly bounded by
\[
\bigl|v_{i,C_t}^{x^*}(b)-h_{i,C_t[b]}^*(X_{t+1})\bigr|
\le
(D_\eta-1)+H_\eta.
\]
By \Cref{prop:public-routing}\textnormal{(R3)}, along any play path there are at
most $D_\eta$ strict child moves. Hence
\begin{equation}
\label{eq:psi-bound}
|Z_{i,t}(\varsigma)|
\le
D_\eta\bigl(D_\eta-1+H_\eta\bigr),
\qquad
|\Psi_{i,t}(\varsigma)|
\le
\Psi_\eta^{\max}:=
H_\eta+D_\eta\bigl(D_\eta-1+H_\eta\bigr)
\end{equation}
for all $t$, all players $i$, and all routed profiles $\varsigma$. This is the
explicit endpoint control missing from v20.

Define the total stagewise gain term
\begin{equation}
\label{eq:gain-chain-def}
\mathfrak G_{i,T}(\varsigma)
:=
\sum_{t\in\mathcal T_T(\varsigma)} g_{i,C_t}^*
=
\sum_{V\in\mathcal V_T(\varsigma)} g_{i,C(V)}^*\,\ell(V),
\end{equation}
and the total complete-tail length
\[
L_T(\varsigma):=|\mathcal T_T(\varsigma)|
=
\sum_{V\in\mathcal V_T(\varsigma)}\ell(V).
\]

\begin{lemma}[Stagewise potential drift]
\label{lem:visit-bellman}
There exists a constant $K_\eta^{\mathrm{pot}}<\infty$ and a finite constant
$A_\eta^{\mathrm{pot}}\ge 1$ such that for every player $i$ the following holds.

\begin{enumerate}[label=\textnormal{(\alph*)},leftmargin=2.8em]
\item For the compliant profile $\sigma_{\eta,\delta,M_*}$,
\begin{equation}
\label{eq:stagewise-potential-compliant}
\E\left[
\sum_{t\in\mathcal T_T(\sigma_{\eta,\delta,M_*})}
u_i^t(\sigma_{\eta,\delta,M_*})
\right]
\ge
\E\bigl[\mathfrak G_{i,T}(\sigma_{\eta,\delta,M_*})\bigr]
-
C_\eta\delta\,\E\bigl[L_T(\sigma_{\eta,\delta,M_*})\bigr]
-
K_\eta^{\mathrm{pot}}.
\end{equation}

\item For every unilateral deviation $\tau_i$ and
$\varsigma^\tau:=(\tau_i,(\sigma_{\eta,\delta,M_*})_{-i})$,
\begin{equation}
\label{eq:stagewise-potential-deviator}
\E\left[
\sum_{t\in\mathcal T_T(\varsigma^\tau)}
u_i^t(\varsigma^\tau)
\right]
\le
\E\bigl[\mathfrak G_{i,T}(\varsigma^\tau)\bigr]
+
A_\eta^{\mathrm{pot}}\eta\,
\E\bigl[L_T(\varsigma^\tau)\bigr]
+
K_\eta^{\mathrm{pot}}.
\end{equation}
\end{enumerate}
\end{lemma}

\begin{proof}
If $\mathcal T_T(\varsigma)=\varnothing$, then the sums in both displays are zero,
so both claims are immediate after enlarging $K_\eta^{\mathrm{pot}}$ once.
Assume henceforth that the complete tail set is nonempty. Because the frozen
schedule has only a finite prelude before the tail cap is reached and then runs
through contiguous tail visits until the final incomplete fragment, there exist
random times
\[
\underline t_T(\varsigma)\le \overline t_T(\varsigma)
\]
such that
\[
\mathcal T_T(\varsigma)
=
\{\underline t_T(\varsigma),\underline t_T(\varsigma)+1,\dots,
\overline t_T(\varsigma)\}.
\]

Fix a stage $t\in\mathcal T_T(\varsigma)$ lying inside a complete tail visit in
node $C=C_t$, with current phase-state $\xi=X_t$. For any action $a_i\in A_i(\xi)$,
the definition of $Z$ gives
\begin{align}
\label{eq:psi-one-step-expansion}
\E\bigl[\Psi_{i,t+1}(\varsigma)\mid \cF_t,\ a_i^t=a_i\bigr]
&=
Z_{i,t}(\varsigma)
+
\sum_{\xi'}Q_{i,C}^{\sigma_{-i,C}^*}(\xi'\mid\xi,a_i)\,h_{i,C}^*(\xi')
\nonumber\\
&\qquad
+
\sum_{b\in B_C}\Lambda_{i,C}^{\sigma_{-i,C}^*}(b\mid\xi,a_i)\,
v_{i,C}^{x^*}(b).
\end{align}
Indeed, if the current visit stays in node $C$, then $Z$ does not move and the
next potential contribution is $h_{i,C}^*(X_{t+1})$; if stage $t$ exits through
$b$, then the increment in \eqref{eq:psi-stock-update} cancels the child-node
bias and leaves exactly the Bellman continuation credit $v_{i,C}^{x^*}(b)$.

For a unilateral deviator $\varsigma^\tau$, the nondeviators still generate the
frozen one-step coefficients of the Bellman problem. Combining
\eqref{eq:psi-one-step-expansion} with the Bellman inequalities in
\eqref{eq:Bellman-ineq-def} yields
\[
\E\bigl[u_i^t(\varsigma^\tau)\mid \cF_t\bigr]
\le
g_{i,C_t}^*
+
\Psi_{i,t}(\varsigma^\tau)
-
\E\bigl[\Psi_{i,t+1}(\varsigma^\tau)\mid \cF_t\bigr]
+\eta
\qquad (t\in\mathcal T_T(\varsigma^\tau)).
\]
Summing over the interval $\mathcal T_T(\varsigma^\tau)$ and using
\eqref{eq:psi-bound} gives
\[
\E\left[
\sum_{t\in\mathcal T_T(\varsigma^\tau)}
u_i^t(\varsigma^\tau)
\right]
\le
\E\bigl[\mathfrak G_{i,T}(\varsigma^\tau)\bigr]
+
\eta\,\E\bigl[L_T(\varsigma^\tau)\bigr]
+
2\Psi_\eta^{\max}.
\]
This is \eqref{eq:stagewise-potential-deviator} after taking
$A_\eta^{\mathrm{pot}}\ge 1$ and $K_\eta^{\mathrm{pot}}\ge 2\Psi_\eta^{\max}$.

Under the compliant profile, every action in the support of
$\sigma_{i,C_t}^*(\cdot\mid X_t)$ attains equality in the Bellman system, so the
previous display is an equality without the $+\eta$ term:
\[
\E\bigl[u_i^t(\sigma_{\eta,\delta,M_*})\mid \cF_t\bigr]
=
g_{i,C_t}^*
+
\Psi_{i,t}(\sigma_{\eta,\delta,M_*})
-
\E\bigl[\Psi_{i,t+1}(\sigma_{\eta,\delta,M_*})\mid \cF_t\bigr].
\]
Summing over $\mathcal T_T(\sigma_{\eta,\delta,M_*})$ and using
\eqref{eq:psi-bound} yields
\[
\E\left[
\sum_{t\in\mathcal T_T(\sigma_{\eta,\delta,M_*})}
u_i^t(\sigma_{\eta,\delta,M_*})
\right]
=
\E\bigl[\mathfrak G_{i,T}(\sigma_{\eta,\delta,M_*})\bigr]
+
\E\bigl[
\Psi_{i,\underline t_T(\sigma_{\eta,\delta,M_*})}
-
\Psi_{i,\overline t_T(\sigma_{\eta,\delta,M_*})+1}
\bigr]
\]
and hence
\[
\E\left[
\sum_{t\in\mathcal T_T(\sigma_{\eta,\delta,M_*})}
u_i^t(\sigma_{\eta,\delta,M_*})
\right]
\ge
\E\bigl[\mathfrak G_{i,T}(\sigma_{\eta,\delta,M_*})\bigr]
-
2\Psi_\eta^{\max}.
\]
Since $C_\eta\delta\,\E[L_T(\sigma_{\eta,\delta,M_*})]\ge 0$, this implies
\eqref{eq:stagewise-potential-compliant} after taking
$K_\eta^{\mathrm{pot}}\ge 2\Psi_\eta^{\max}$.
\end{proof}

\begin{remark}[What changed from v20]
The stagewise telescope variable is not the raw bias $\Phi_{i,t}$ but the
stock-augmented potential $\Psi_{i,t}=h_{i,C_t}^*(X_t)+Z_{i,t}$. The explicit
stock $Z$ repairs the boundary mismatch between a parent-node Bellman credit
$v_{i,C}^{x^*}(b)$ and the successor-node entry bias. Version~26 keeps this
process as the correct local Bellman-transport device across node exits, but the
final global comparison is now taken from the optimal bias of the induced MDP
itself rather than from a routed sum of local gains.
\end{remark}

\begin{lemma}[Global comparison in the induced MDP]
\label{lem:global-mdp-comparison}
Fix a player $i$, the realized profile $\sigma_{\eta,\delta,M_*}$, an initial
state $s$, a horizon $T$, and a unilateral deviation $\tau_i$. Write
\[
\varsigma^\sigma:=\sigma_{\eta,\delta,M_*},
\qquad
\varsigma^\tau:=(\tau_i,(\sigma_{\eta,\delta,M_*})_{-i}).
\]
Then there exist finite constants
$A_\eta^{\mathrm{mdp}},K_\eta^{\mathrm{mdp}}<\infty$ such that
\begin{equation}
\label{eq:global-mdp-comparison}
\gamma_i^T\bigl(s,\varsigma^\tau\bigr)
-
\gamma_i^T\bigl(s,\varsigma^\sigma\bigr)
\le
A_\eta^{\mathrm{mdp}}\eta + C_\eta\delta + \frac{K_\eta^{\mathrm{mdp}}}{T}.
\end{equation}
\end{lemma}

\begin{proof}
Apply \Cref{prop:mdp-global-comparison} to the compliant profile
$\varsigma^\sigma$ and the deviation profile $\varsigma^\tau$. That proposition
compares both profiles through the global optimal gain $g_i^*$ of the induced
unichain MDP and yields
\[
\gamma_i^T\bigl(s,\varsigma^\tau\bigr)
-
\gamma_i^T\bigl(s,\varsigma^\sigma\bigr)
\le
A_\eta^{\mathrm{opt}}\eta + C_\eta\delta + \frac{2K_\eta^{\mathrm{mdp}}}{T}.
\]
After enlarging $A_\eta^{\mathrm{mdp}}$ and $K_\eta^{\mathrm{mdp}}$ to dominate the
right-hand constants, this is exactly \eqref{eq:global-mdp-comparison}.
\end{proof}

\begin{remark}[What v26 no longer compares]
The local gains $g_{i,C}^*$ remain the local Bellman data used to define the
fixed point and the compliant kernels. They are not compared across compliant
and deviating node sequences, no gain-near-constancy statement is invoked, and
no identity of the form
\[
\sum_t g_{i,C_t}^*=g_i^* T
\]
enters the proof. The final benchmark is the global optimal gain $g_i^*$ of the
induced one-player MDP.
\end{remark}

Fix once and for all a constant
\[
A_\eta\ge 1+A_\eta^{\mathrm{mdp}}.
\]


\subsection{Proof of the main theorem}


\begin{proof}[Proof of \Cref{thm:main}]
Fix $\eps>0$. Choose parameters in the order
\[
\eta,\qquad \delta,\qquad M_*,\qquad T_0.
\]

\smallskip
\noindent\textbf{Step 1: fixed point and MDP error scale.}
Choose $\eta>0$ so small that
\[
A_\eta\eta\le \frac{\eps}{3}.
\]
By \Cref{thm:kakutani-global}, there exists a fixed point $x^*\in\BR_\eta(x^*)$.

\smallskip
\noindent\textbf{Step 2: realization scale.}
Let $C_\eta$ dominate the finitely many $\delta$-loss constants appearing in
\Cref{thm:realization,lem:global-mdp-comparison}. Choose $\delta>0$ such that
\[
C_\eta\delta\le \frac{\eps}{3}.
\]
Since the tree is finite, there is a common cap threshold
\[
L_\eta(\delta):=\max_{C\in\cT_\eta}L_C(\delta)
\]
from \Cref{thm:local-controller}. Freeze a superexponential schedule at a level
$R$ so large that the associated tail cap $M_*=M_R$ satisfies
$M_*\ge L_\eta(\delta)$. Let
\[
\sigma^\eps:=\sigma_{\eta,\delta,M_*}
\]
be the realized profile from \Cref{thm:realization}.

\smallskip
\noindent\textbf{Step 3: horizon threshold.}
Let $K_\eta<\infty$ dominate the remainder constant from
\Cref{lem:global-mdp-comparison}. Choose $T_0$ so large that
\[
\frac{K_\eta}{T_0}\le \frac{\eps}{3}.
\]

Now fix $T\ge T_0$, an initial state $s$, a player $i$, and a unilateral
deviation $\tau_i$. By \Cref{lem:global-mdp-comparison},
\begin{equation}
\label{eq:global-mdp-comparison-proof}
\gamma_i^T\bigl(s,(\tau_i,\sigma^\eps_{-i})\bigr)
-
\gamma_i^T(s,\sigma^\eps)
\le
A_\eta\eta + C_\eta\delta + \frac{K_\eta}{T}.
\end{equation}
By construction, the right-hand side is at most
\[
\frac{\eps}{3}+\frac{\eps}{3}+\frac{\eps}{3}=\eps.
\]
Therefore
\[
\gamma_i^T(s,\sigma^\eps)
\ge
\gamma_i^T\bigl(s,(\tau_i,\sigma^\eps_{-i})\bigr)-\eps.
\]
Since $T\ge T_0$, $s$, $i$, and $\tau_i$ were arbitrary, $\sigma^\eps$ is a
public finite-memory uniform $\eps$-equilibrium.
\end{proof}



\section{Discussion}

The full Neyman conjecture remains open. What the present paper proves is the
downstream implication from a finite universal $\eta$-tree plus its first-exit
controller realization to a public finite-memory uniform $\eps$-equilibrium,
under condition~\textnormal{(a)} and the global induced-MDP comparison from
\Cref{lem:global-mdp-comparison}.

The conceptual point of v26 is that the final benchmark is global from the start.
The identity
\[
g_{i,C}^*
=
\alpha_{i,C}^*
+
\sum_{b\in B_C}\beta_C^*(b)\,v_{i,C}^{x^*}(b)
\]
belongs to the compliant node data and remains useful only there. It is part of
the construction of the fixed point and of the compliant local kernels, but it
is not compared across different routed node sequences and it is not used as a
global scalar benchmark.

Once player $i$ deviates, the other players continue to play
$(\sigma_{\eta,\delta,M_*})_{-i}$, so player $i$ faces a finite one-player MDP
with fixed environment. The final regret estimate then comes from standard
average-reward MDP theory. The global optimal gain $g_i^*$ and optimal bias
$h_i^*$ bound every unilateral deviation from above by
$g_i^*+O(1/T)$. The compliant stationary policy selected by the fixed point is
$O(\eta+\delta)$-optimal in that same induced MDP, and its own bias gives the
matching lower bound
$g_i^*-O(\eta+\delta+1/T)$. Subtracting the two estimates yields the uniform
regret bound.

This is exactly the repair required by Pass~114: the local tree machinery is
retained intact, but the failed local-to-global bridge is removed. No routed
comparison of local gains is needed in the final step.
\appendix




\section{Finite average-reward MDP facts and the global induced-MDP comparison}
\label{app:mdp}

This appendix records the one-player facts used in the final comparison.
All statements are standard; see Puterman's monograph \cite{Puterman1994} for
background. The point specific to v26 is that once player $i$ deviates and the
other players keep the realized profile, player $i$ faces a finite induced MDP,
and the regret estimate is obtained directly from the global gain-bias theory of
that MDP rather than from routed comparisons of local gains.

Consider a finite average-reward MDP with state space $X$, action sets $A(x)$,
bounded reward function $r(x,a)\in[0,\overline r]$ for some $\overline r<\infty$,
and transition kernel $Q(\cdot\mid x,a)$.

\begin{proposition}[Primal-dual formulation]
\label{prop:mdp-duality}
The optimal average reward is the common value of the primal and dual linear
programs
\begin{align*}
\textnormal{maximize }&\sum_{x\in X}\sum_{a\in A(x)}\nu(x,a)r(x,a)
\\
\textnormal{subject to }&\nu\ge 0,
\qquad \sum_{x,a}\nu(x,a)=1,
\\
&\sum_{a\in A(y)}\nu(y,a)=\sum_{x,a}\nu(x,a)Q(y\mid x,a)
\qquad (y\in X),
\end{align*}
and
\begin{align*}
\textnormal{minimize }&g
\\
\textnormal{subject to }&g+h(x)\ge r(x,a)+\sum_y Q(y\mid x,a)h(y)
\qquad (x\in X,\ a\in A(x)).
\end{align*}
There exists an optimal stationary mixed policy and an optimal dual pair $(g,h)$.
Moreover, after the linear normalization $h(x^\circ)=0$ at any fixed reference
state $x^\circ\in X$, one may choose $h$ in a compact box.
\end{proposition}

\begin{proof}
This is the classical average-reward LP formulation for finite MDPs. Feasibility is
immediate. The primal feasible set is a nonempty compact polytope, so an optimum
exists. The dual feasible set is nonempty because choosing $g=\overline r$ and
$h\equiv 0$ is feasible. Finite-dimensional LP duality gives equality of primal
and dual
optima and existence of optimal primal and dual solutions. The occupation measure
of an optimal primal solution induces an optimal stationary mixed policy. The
normalization $h(x^\circ)=0$ is obtained by subtracting the constant $h(x^\circ)$.
\end{proof}

\begin{proposition}[Uniqueness of the normalized bias in the communicating case]
\label{prop:bias-unique}
If the finite average-reward MDP is communicating, then any two optimal bias
vectors differ by a constant. Consequently, after fixing one reference coordinate
$h(x^\circ)=0$, the optimal bias is unique.
\end{proposition}

\begin{proof}
This is the standard communicating-MDP relative-value theorem; see
\cite[Chapter~8]{Puterman1994}. In the communicating case the optimal average
reward is state-independent and every optimal bias vector is unique up to an
additive constant. Fixing one coordinate removes that one-dimensional freedom.
\end{proof}

\begin{proposition}[Unichain optimality equation and finite-horizon comparison]
\label{prop:unichain-horizon}
Assume that the finite average-reward MDP is unichain. Then there exist a scalar
optimal gain $g^*$ and a bias $h^*:X\to\R$ such that
\[
g^*+h^*(x)=\max_{a\in A(x)}\Bigl(r(x,a)+\sum_y Q(y\mid x,a)h^*(y)\Bigr)
\qquad (x\in X).
\]
For every admissible policy $\pi$ (stationary or history-dependent), every
initial state $x\in X$, and every horizon $T\ge 1$,
\[
\frac1T\,\E_x^\pi\!\left[\sum_{t=1}^T r(X_t,A_t)\right]
\le
g^*+\frac{\spn(h^*)}{T}.
\]
If $\pi$ is stationary and $(g^\pi,h^\pi)$ is any gain-bias pair satisfying
\[
g^\pi+h^\pi(x)=r^\pi(x)+\sum_y Q^\pi(y\mid x)h^\pi(y)
\qquad (x\in X),
\]
then for every initial state $x$ and every horizon $T\ge 1$,
\[
\frac1T\,\E_x^\pi\!\left[\sum_{t=1}^T r(X_t,A_t)\right]
=
g^\pi+\frac{h^\pi(x)-\E_x^\pi[h^\pi(X_{T+1})]}{T},
\]
and in particular
\[
g^\pi-\frac{\spn(h^\pi)}{T}
\le
\frac1T\,\E_x^\pi\!\left[\sum_{t=1}^T r(X_t,A_t)\right]
\le
g^\pi+\frac{\spn(h^\pi)}{T}.
\]
\end{proposition}

\begin{proof}
Existence of $(g^*,h^*)$ is the standard unichain average-reward optimality
equation; see \cite[Chapter~8]{Puterman1994}. Its maximizing form implies the
Bellman inequality
\[
r(x,a)-g^*
\le
h^*(x)-\sum_y Q(y\mid x,a)h^*(y)
\qquad (x\in X,\ a\in A(x)).
\]
Let $(X_t,A_t)$ evolve under an arbitrary admissible policy $\pi$ from initial
state $x$. Define
\[
M_t:=h^*(X_t)-\sum_{k=1}^{t-1}\bigl(r(X_k,A_k)-g^*\bigr)
\qquad (t\ge 1).
\]
The displayed Bellman inequality gives
\[
\E_x^\pi[M_{t+1}\mid \cF_t]\le M_t,
\]
so $(M_t)_{t\ge 1}$ is a supermartingale. Therefore
\[
\E_x^\pi[M_{T+1}]\le M_1=h^*(x),
\]
which rearranges to
\[
\E_x^\pi\!\left[\sum_{t=1}^T r(X_t,A_t)\right]
\le
Tg^*+h^*(x)-\E_x^\pi[h^*(X_{T+1})]
\le
Tg^*+\spn(h^*).
\]
Dividing by $T$ gives the first bound.

Now let $\pi$ be stationary. Its Poisson equation implies that the same process
built from $(g^\pi,h^\pi)$ is a martingale, hence
\[
\E_x^\pi\!\left[\sum_{t=1}^T r(X_t,A_t)\right]
=
Tg^\pi+h^\pi(x)-\E_x^\pi[h^\pi(X_{T+1})].
\]
Dividing by $T$ gives the exact identity, and the span bounds are immediate.
\end{proof}

\begin{remark}
For every fixed finite unichain MDP, the spans $\spn(h^*)$ and $\spn(h^\pi)$ are
finite. One may bound them explicitly in terms of the MDP diameter and the reward
range; see again \cite[Chapter~8]{Puterman1994}. In the main text we use only
finiteness and absorb the resulting constants into $K_\eta$.
\end{remark}


\begin{proposition}[The unilateral-control problem is a finite unichain MDP]
\label{prop:induced-routed-mdp}
Fix a player $i$, a fixed point $x^*\in\BR_\eta(x^*)$, and the realized profile
$\sigma_{\eta,\delta,M_*}$. If the other players keep
$(\sigma_{\eta,\delta,M_*})_{-i}$, then player $i$ faces a finite average-reward
MDP on a finite public-memory state space. Its state records the current
stochastic-game state together with the public memory of the realized profile,
including the current routed node, the current local controller state, and the
finite memory of the frozen cap schedule. The stage reward is the expected stage
payoff against the frozen opponents, and the transition kernel is the induced law
of the next public-memory state. Every unilateral deviation strategy $\tau_i$
determines a policy in this MDP, and the compliant behavior of player $i$ is the
stationary policy induced by the local kernels $\sigma_{i,C}^*$. Under
condition~\textnormal{(a)}, this MDP is unichain.
\end{proposition}

\begin{proof}
By \Cref{prop:public-routing}\textnormal{(R4)}, the realized profile uses only
finite public memory. Once the nondeviators' mixed actions are frozen, the law of
the next augmented public-memory state depends only on the current augmented
public-memory state and on player $i$'s current action. The expected stage payoff
at that state is therefore a bounded function of the same data. This is exactly a
finite average-reward MDP.

Any unilateral deviation strategy induces a possibly history-dependent policy of
this MDP, and conversely the compliant local kernels $\sigma_{i,C}^*$ define a
stationary policy on the augmented state space. Finally, condition~\textnormal{(a)}
prevents a unilateral deviation from splitting the stabilized tail dynamics into
several communicating classes: the finite prelude states are transient, while the
tail regime remains irreducible under every stationary policy of player $i$.
Hence the induced MDP has a single recurrent class and is unichain.
\end{proof}

\begin{proposition}[Global comparison through the optimal gain of the induced MDP]
\label{prop:mdp-global-comparison}
Fix a player $i$, an initial state $s$, a horizon $T$, a fixed point
$x^*\in\BR_\eta(x^*)$, the compliant profile
$\varsigma^\sigma:=\sigma_{\eta,\delta,M_*}$, and a unilateral deviation
$\varsigma^\tau:=(\tau_i,(\sigma_{\eta,\delta,M_*})_{-i})$. Let $g_i^*$ denote the
optimal gain of the induced unichain MDP from \Cref{prop:induced-routed-mdp}.
Then there exist finite constants
$A_\eta^{\mathrm{opt}},K_\eta^{\mathrm{mdp}}<\infty$ such that for every horizon
$T$,
\begin{align*}
\gamma_i^T\bigl(s,\varsigma^\sigma\bigr)
&\ge
g_i^* - A_\eta^{\mathrm{opt}}\eta - C_\eta\delta - \frac{K_\eta^{\mathrm{mdp}}}{T},
\\
\gamma_i^T\bigl(s,\varsigma^\tau\bigr)
&\le
g_i^* + \frac{K_\eta^{\mathrm{mdp}}}{T}.
\end{align*}
Consequently,
\[
\gamma_i^T\bigl(s,\varsigma^\tau\bigr)
-
\gamma_i^T\bigl(s,\varsigma^\sigma\bigr)
\le
A_\eta^{\mathrm{opt}}\eta + C_\eta\delta + \frac{2K_\eta^{\mathrm{mdp}}}{T}.
\]
In particular, the compliant and deviating profiles are compared through the
single global optimal gain $g_i^*$ of the induced MDP, not through routed
comparisons of the local gains $(g_{i,C}^*)_{C\in\cT_\eta}$.
\end{proposition}

\begin{proof}
By \Cref{prop:induced-routed-mdp}, once the nondeviators are frozen at
$(\sigma_{\eta,\delta,M_*})_{-i}$, player $i$ faces a finite unichain
average-reward MDP. Let $h_i^*$ be an optimal bias for that MDP, and let
$(g_i^\sigma,\widehat h_i^\sigma)$ be a gain-bias pair for the stationary policy
induced by the compliant profile $\varsigma^\sigma$.

The only nonstandard input is the retained output of
Sections~\ref{sec:global}--\ref{sec:realization}: the fixed-point construction,
together with the realized first-exit controllers, produces a stationary policy
whose average reward is $O(\eta+\delta)$-optimal in the induced MDP. Thus
\begin{equation}
\label{eq:compliant-gain-eta-optimal}
g_i^\sigma
\ge
g_i^* - A_\eta^{\mathrm{opt}}\eta - C_\eta\delta.
\end{equation}

Apply \Cref{prop:unichain-horizon} to the optimal pair $(g_i^*,h_i^*)$ and the
possibly history-dependent policy generated by $\varsigma^\tau$. Since the
initial stochastic-game state $s$ and the initial public memory determine one
initial MDP state, this yields
\[
\gamma_i^T\bigl(s,\varsigma^\tau\bigr)
\le
g_i^*+\frac{\spn(h_i^*)}{T}.
\]
Apply the stationary-policy part of \Cref{prop:unichain-horizon} to the
compliant stationary policy $\varsigma^\sigma$:
\[
\gamma_i^T\bigl(s,\varsigma^\sigma\bigr)
\ge
g_i^\sigma-\frac{\spn(\widehat h_i^\sigma)}{T}.
\]
Combining this with \eqref{eq:compliant-gain-eta-optimal} gives
\[
\gamma_i^T\bigl(s,\varsigma^\sigma\bigr)
\ge
g_i^* - A_\eta^{\mathrm{opt}}\eta - C_\eta\delta - \frac{\spn(\widehat h_i^\sigma)}{T}.
\]
Now choose $K_\eta^{\mathrm{mdp}}$ large enough to dominate both
$\spn(h_i^*)$ and $\spn(\widehat h_i^\sigma)$. The two displayed inequalities and
the regret bound follow immediately.
\end{proof}



\section{Summing stagewise drift inequalities}
\label{app:aggregation}

\begin{lemma}
\label{lem:aggregation}
Let $(\cF_t)_{t=1}^{T+1}$ be a filtration, let $(Y_t)_{t=1}^{T+1}$ be an adapted
process with $|Y_t|\le H$ almost surely for every $t$, and let
$(\Delta_t)_{t=1}^T$ and $(a_t)_{t=1}^T$ be integrable. If
\[
\E[\Delta_t\mid \cF_t]
\le
a_t + Y_t - \E[Y_{t+1}\mid \cF_t]
\qquad (1\le t\le T),
\]
then
\[
\E\left[\sum_{t=1}^T \Delta_t\right]
\le
\E\left[\sum_{t=1}^T a_t\right] + 2H.
\]
The same conclusion holds with the reverse inequality throughout.
\end{lemma}

\begin{proof}
Taking expectations and summing over $t$ gives
\[
\E\left[\sum_{t=1}^T \Delta_t\right]
\le
\E\left[\sum_{t=1}^T a_t\right]
+
\E[Y_1-Y_{T+1}],
\]
and the last term is bounded in absolute value by $2H$.
\end{proof}

\begin{thebibliography}{99}


\bibitem{BewleyKohlberg1976}
T.~Bewley and E.~Kohlberg.
\newblock The asymptotic theory of stochastic games.
\newblock \emph{Mathematics of Operations Research}, 1(3):197--208, 1976.

\bibitem{FleschThuijsmanVrieze1997}
J.~Flesch, F.~Thuijsman, and K.~Vrieze.
\newblock Cyclic {M}arkov equilibria in stochastic games.
\newblock \emph{International Journal of Game Theory}, 26:303--314, 1997.

\bibitem{Kakutani1941}
S.~Kakutani.
\newblock A generalization of {B}rouwer's fixed point theorem.
\newblock \emph{Duke Mathematical Journal}, 8(3):457--459, 1941.

\bibitem{MertensNeyman1981}
J.-F.~Mertens and A.~Neyman.
\newblock Stochastic games.
\newblock \emph{International Journal of Game Theory}, 10(2):53--66, 1981.

\bibitem{MertensSorinZamir2003}
J.-F.~Mertens, S.~Sorin, and S.~Zamir.
\newblock Stochastic games.
\newblock In R.~Aumann and S.~Hart, editors, \emph{Handbook of Game Theory with
Economic Applications}, volume~3, chapter~47, pages 1811--1832. Elsevier,
Amsterdam, 2003.

\bibitem{Neyman2003}
A.~Neyman.
\newblock From {M}arkov chains to stochastic games.
\newblock In A.~Neyman and S.~Sorin, editors, \emph{Stochastic Games and
Applications}, NATO Science Series C, volume~570, pages 9--25. Kluwer Academic
Publishers, Dordrecht, 2003.

\bibitem{Nowak2003}
A.~S.~Nowak.
\newblock N-person stochastic games: Extensions of the finite state space case and
correlation.
\newblock In A.~Neyman and S.~Sorin, editors, \emph{Stochastic Games and
Applications}, NATO Science Series C, volume~570, pages 93--106. Kluwer Academic
Publishers, Dordrecht, 2003.

\bibitem{Puterman1994}
M.~L.~Puterman.
\newblock \emph{Markov Decision Processes: Discrete Stochastic Dynamic
Programming}.
\newblock John Wiley \& Sons, New York, 1994.

\bibitem{Shapley1953}
L.~S.~Shapley.
\newblock Stochastic games.
\newblock \emph{Proceedings of the National Academy of Sciences of the United
States of America}, 39(10):1095--1100, 1953.

\bibitem{Solan1999}
E.~Solan.
\newblock Three-player absorbing games.
\newblock \emph{Mathematics of Operations Research}, 24(3):669--698, 1999.

\bibitem{Solan2003}
E.~Solan.
\newblock Uniform equilibrium: More than two players.
\newblock In A.~Neyman and S.~Sorin, editors, \emph{Stochastic Games and
Applications}, NATO Science Series C, volume~570, pages 285--294. Kluwer Academic
Publishers, Dordrecht, 2003.

\bibitem{SolanVieille2002}
E.~Solan and N.~Vieille.
\newblock Correlated equilibrium in stochastic games.
\newblock \emph{Games and Economic Behavior}, 38(2):362--399, 2002.

\bibitem{Vieille2000I}
N.~Vieille.
\newblock Equilibrium in two-person stochastic games {I}: A reduction.
\newblock \emph{Israel Journal of Mathematics}, 119:55--91, 2000.

\bibitem{Vieille2000II}
N.~Vieille.
\newblock Equilibrium in two-person stochastic games {II}: The case of recursive
games.
\newblock \emph{Israel Journal of Mathematics}, 119:93--126, 2000.

\bibitem{Vieille2002}
N.~Vieille.
\newblock Stochastic games: Recent results.
\newblock In R.~Aumann and S.~Hart, editors, \emph{Handbook of Game Theory with
Economic Applications}, volume~3, chapter~48, pages 1833--1850. Elsevier,
Amsterdam, 2002.

\end{thebibliography}

\end{document}
